\begin{textsample}{NEG Dim 1 – 2019\_c6 – Score -5.00 – 2019\_c6/enade\_2019\_eng\_flo\_3a2b53498b.txt}  \label{ex:f1_neg_098}
CAPÍTULO 6 ANÁLISE TÉCNICA DA PROVA Este capítulo tem por objetivo apresentar o desempenho dos estudantes concluintes de Engenharia Florestal no Enade/2019.
Para isso, foram calculadas as estatísticas básicas da prova em seu todo ( seção 6.1.1 ), bem como as estatísticas relacionadas aos Componentes de Formação Geral ( seção 6.1.2 ) e de Conhecimento Específico da Área ( seção 6.1.3 ).
Dadas as suas características, foram analisadas, em separado, as questões objetivas ( seção 6.2 ) e as questões discursivas ( seção 6.3 ).
Tomando -se como base as duas questões discursivas do Componente de Formação Geral, nas seções 6.3.1.5 a 6.3.1.8, são apresentados comentários sobre a correção das respostas em relação à Língua Portuguesa.
Para as questões objetivas, foram disponibilizados os índices de facilidade e de discriminação Ponto-Bisserial, também em separado, para os Componentes de Formação Geral ( seção 6.2.1 ) e de Conhecimento Específico ( seção 6.2.2 ).
De \textbf{cada} componente, uma das questões foi escolhida para exemplificar a análise gráfica, relacionando as alternativas escolhidas pelos estudantes ( inclusive o gabarito ) com o número de acertos no componente.
No Anexo I, apresenta -se a íntegra da análise gráfica para todas as questões objetivas.
Para \textbf{cada} uma das questões discursivas, os conteúdos dos tipos mais comuns de respostas dos estudantes são apresentados e comparados com o padrão de resposta esperado ( ver Anexo VIII com o padrão de respostas ).
Nas tabelas, constam as seguintes estatísticas das notas¹⁹ : média do desempenho na prova, erro padrão da média, desvio padrão, nota mínima, mediana e nota máxima para \textbf{cada} um de seus componentes.
Tais estatísticas contemplam o total de estudantes concluintes da Área de Engenharia Florestal inscritos e presentes à prova do Enade/2019, tendo em vista agregações, ou por Grandes Regiões e o país como um todo, ou por Categoria Administrativa e Organização Acadêmica da IES.
Em relação aos gráficos de distribuição de notas, o intervalo considerado foi de 10 unidades, aberto à esquerda e fechado à direita, com exceção do primeiro intervalo, [ 0 ; 10 ], fechado em ambos os extremos.
Para os gráficos de distribuição das notas das questões discursivas, foram consideradas mais duas categorias : questão em branco²⁰ e nota zero. ¹⁹ Uma definição dessas estatísticas pode ser encontrada no Glossário. ²⁰ Nesse grupo \textbf{estão} incluídas também as respostas classificadas como nulas ou desconsideradas. 6.1 ESTATÍSTICAS BÁSICAS DA PROVA Esta seção apresenta estatísticas selecionadas e histogramas da nota geral ( 6.1.1 ) e de \textbf{cada} componente : Formação Geral ( 6.1.2 ) e Conhecimento Específico ( 6.1.3 ).
São, também, apresentadas estatísticas selecionadas de subpopulações, caracterizadas por Grande Região, Categoria Administrativa e Organização Acadêmica. 6.1.1 Estatísticas Básicas Gerais Na Tabela 6.1, são apresentadas as Estatísticas Básicas da prova, por Grande Região, dos estudantes concluintes de Engenharia Florestal.
A Média das notas da prova, como um todo ( nas seções seguintes serão analisados os Componentes de Formação Geral e de Conhecimento Específico ), foi 44,1, sendo que os estudantes da região Norte obtiveram a Média mais baixa ( 38,8 ), e os da região Sudeste obtiveram a Média mais alta ( 48,9 ).
As demais Médias foram : 42,8, na região Nordeste, 46,2, na região Sul e 44,8, na região Centro-Oeste.
O Desvio padrão para o Brasil, como um todo, foi 12,2, sendo este o maior Desvio padrão encontrado na região Sul ( 12,2 ), e o menor, na região Norte ( 11,1 ), indicando uma dispersão um pouco menor das notas dessa última região.
A região que obteve a maior nota Máxima foi a Sudeste ( 85,9 ), ao passo que a região que atingiu a menor nota Máxima foi a Norte ( 70,2 ).
A Mediana do Brasil, como um todo, foi 44,0, sendo a maior Mediana obtida na região Sudeste ( 49,6 ), e a menor obtida na região Norte ( 38,3 ).
A nota Mínima do Brasil foi obtida na região Norte, onde esta foi 7,6.
A maior nota Mínima, entre as cinco Grandes Regiões, foi 15,4, na região Centro-Oeste.
As demais regiões obtiveram notas Mínima iguais à 12,4 ( Nordeste e Sudeste ) e 12,9 ( Sul ).
Considerando -se as notas segundo Grande Região, observa -se que existem diferenças estatisticamente significativas ao nível de 95 \% entre as Médias, em todas as regiões, à exceção da diferença entre a região Centro-Oeste e as regiões Nordeste e Sul.
O comportamento das notas dos estudantes de todo o Brasil pode ser observado no Gráfico 6.1, que apresenta um histograma com a distribuição das mesmas : uma distribuição unimodal com a moda no intervalo ( 40 ; 50 ].
Na Tabela 6.2, são apresentadas informações referentes à Média da nota final, desagregadas por Categoria Administrativa e por Organização Acadêmica.
A Média das notas da prova, como um todo, foi 44,1.
Em relação a Categoria Administrativa, os estudantes das IES Públicas obtiveram Média mais alta ( 45,0 ) que a Média nacional, e os das IES Privadas obtiveram Média mais baixa ( 35,8 ).
Observa -se que existe diferença estatisticamente significativa entre as Médias das notas das IES Públicas e as das Privadas.
A diferença entre as Médias das regiões Sudeste e Norte ( 10,1 ), a maior e a menor Média, é superior à diferença entre a das IES Públicas e Privadas ( 9,2 ), caracterizando -se uma menor diversidade administrativa do que regional.
O Desvio padrão para as IES Públicas ( 12,0 ) e Privadas ( 11,4 ), ambos inferiores ao do Brasil, como um todo ( 12,2 ), indicando uma dispersão um pouco menor das notas por Categoria Administrativa do que entre os cursos em \textbf{cada} uma delas.
No tocante a Organização Acadêmica, as Universidades e os CEFET/IF obtiveram Média mais alta que a nacional ( 44,9 e 48,0, respectivamente ).
A Média dos Centros Universitários e a das Faculdades foram menores do que a nacional ( 33,0 e 33,9, respectivamente ).
Constata -se que existe diferença estatisticamente significativa, ao nível de 95 \%, entre a Média das Faculdades e a Média das Universidades e os CEFET/IF. 6.1.2 Estatísticas Básicas no Componente de Formação Geral Na Tabela 6.3, são apresentadas as Estatísticas Básicas em relação ao componente da prova que avalia a Formação Geral dos estudantes concluintes.
A nota de Formação Geral é obtida levando -se em conta os acertos nas oito questões objetivas e a nota média obtida nas questões discursivas 1 e 2, sendo a nota de \textbf{cada} questão discursiva obtida levando -se em conta 80 \% da nota relativa ao conteúdo e 20 \% da nota relativa ao desempenho linguístico.
Os estudantes de todo o Brasil obtiveram desempenho médio de 47,2.
Quanto à variabilidade, o Desvio padrão das notas dos estudantes do Brasil, como um todo, foi 16,4.
A maior Média foi obtida na região Sudeste ( 52,2 ), e a menor, na região Norte ( 42,5 ).
As demais Médias foram : 47,2, na região Nordeste, 47,1, na região Sul e 47,8, na região Centro-Oeste.
Já o maior Desvio padrão foi obtido na região Nordeste ( 16,6 ) e o menor, na região Sudeste ( 15,2 ).
A maior nota no Componente de Formação Geral da prova do Enade/2019 foi 97,8, obtida por, pelo menos, um estudante na região Centro-Oeste.
A menor nota Máxima foi obtida na região Norte ( 84,6 ).
A Mediana do Brasil, como um todo, foi 46,9, a mesma da região Centro-Oeste, sendo a menor Mediana encontrada na região Norte ( 41,7 ), e a maior, na região Sudeste ( 54,1 ).
A nota Mínima nessa parte foi zero somente na região Sul.
Considerando -se as notas, segundo Grande Região, observa -se que existem diferenças estatisticamente significativas entre as Médias das notas do Componente de Formação Geral obtidas nas regiões Norte e Sudeste, e as demais regiões.
No Gráfico 6.2, é possível verificar a avaliação do desempenho dos estudantes do Componente de Formação Geral, a partir do histograma da distribuição das notas correspondentes.
A distribuição é unimodal, com moda em ( 40 ; 50 ], intervalo exatamente igual ao modal da distribuição de notas da prova, como um todo ( Gráfico 6.1 ).
Nota -se, ainda, que, no Gráfico 6.2, as notas apresentam maior dispersão do que as no Gráfico 6.1 ( distribuição das notas da prova ), confirmada pela comparação dos desvios padrões : 12,2 para a nota da prova, como um todo, e 16,4 para o Componente de Formação Geral.
Na Tabela 6.4, são apresentadas as informações referentes ao desempenho dos concluintes do Componente de Formação Geral, em diferentes agregações : Categoria Administrativa e Organização Acadêmica.
Observa -se que existe diferença estatisticamente significativa entre as Médias dos tipos de Categoria Administrativa : a maior Média obtida por estudantes de IES Públicas ( 48,3 ) e a menor, pelos de IES Privadas ( 37,0 ).
Nota -se que, também, há diferença estatisticamente significativa entre as três maiores Médias dos quatro tipos de Organização Acadêmica que englobam os CEFET/IF ( 48,5 ), com as Universidades ( 48,1 ) e as Faculdades ( 36,6 ).
Sendo, os Centros Universitários ( 34,4 ), a única Organização Acadêmica que não apresenta diferença estatisticamente significativa com nenhuma outra. 6.1.3 Estatísticas Básicas do Componente de Conhecimento Específico Na Tabela 6.5, são apresentadas as Estatísticas Básicas referentes ao Componente de Conhecimento Específico da área de Engenharia Florestal.
A nota do componente de Conhecimentos Específicos leva em conta as notas da parte objetiva e a média das notas obtidas nas três questões discursivas parte discursiva da Área.
A Média do desempenho dos estudantes do Brasil, como um todo, foi 43,0.
A maior Média foi obtida na região Sudeste ( 47,7 ), e a menor, na região Norte ( 37,6 ).
As demais Médias foram : 41,3, na região Nordeste, 45,9 na região Sul e 43,8, na região Centro-Oeste.
Quanto à variabilidade das notas, o Desvio padrão do Brasil, como um todo, foi 13,3, sendo o maior Desvio padrão observado na região Sudeste ( 13,6 ) e o menor, na região Norte ( 11,9 ).
Os demais desvios foram : 12,7, na região Nordeste, 13,1, na região Sul e 12,2, na região Centro-Oeste.
A Mediana das notas dos estudantes de todo o Brasil foi 42,7.
A maior Mediana ocorreu na região Sudeste ( 47,5 ), e a menor, na região Norte ( 37,4 ).
As demais medianas foram : 41,1 na região Nordeste, 45,5 na região Sul e 42,9 na região Centro-Oeste.
A nota Máxima do Brasil, como um todo, foi 89,9, sendo obtida por, pelo menos, um estudante na região Sudeste.
As demais notas máximas foram : 72,8, na região Norte, 81,8, na região Nordeste, 76,8, na região Sul e 72,3 na região Centro-Oeste.
A nota Mínima do Brasil foi 6,1, sendo essa obtida por, pelo menos, um estudante na região Nordeste.
As demais notas mínimas foram : 6,3, na região Norte, 8,5, na região Sudeste, 8,6, na região Sul e 14,9, na região Centro-Oeste.
Observa -se que existe diferença estatisticamente significativa das notas do Componente de Conhecimento Específico entre a menor Média, obtida na região Norte, e as Médias, das demais regiões, como também entre a segunda menor Média, obtida na região Nordeste, e a Média das regiões Sudeste e Sul.
Assim como nos Gráficos 6.1 e 6.2, no Gráfico 6.3, apresentado a seguir, é possível ser feita uma avaliação do desempenho de concluintes de Engenharia Florestal, em relação ao Componente de Conhecimento Específico, com um histograma da distribuição das notas correspondentes.
Esta também é uma distribuição unimodal, e o grupo modal é o ( 40 ; 50 ], o mesmo do grupo modal da prova, como um todo, e do grupo modal para a Formação Geral.
Na Tabela 6.6, é apresentada uma comparação dos resultados em relação a Categoria Administrativa e a Organização Acadêmica, agora levando -se em conta o desempenho de estudantes do Componente de Conhecimento Específico da prova.
No que se refere a Organização Acadêmica, a maior Média foi a dos CEFET/IF ( 47,8 ), vindo a seguir a das Universidades ( 43,8 ), a das Faculdades ( 33,0 ) e, por fim, a dos Centros Universitários ( 32,5 ).
O maior Desvio padrão, e acima do valor para o Brasil, como um todo, foi o dos Centros Universitários ( 13,5 ).
As Universidades obtiveram a maior nota Máxima ( 89,9 ).
Os Centros Universitários obtiveram nota Máxima 46,9 ; as Faculdades, 62,4, e os CEFET/IF obtiveram nota Máxima 75,6.
As Medianas foram : 43,4 nas Universidades, 30,3 nos Centros Universitários, a menor delas ; 33,2 nas Faculdades e 48,7, nos CEFET/IF, a maior.
Como já citado anteriormente, a nota Mínima do Brasil foi 6,1, sendo essa obtida por, pelo menos, um estudante na Organização Acadêmica Universidades.
As demais notas mínimas foram : 20,2, nos Centros Universitários, 8,6, nas Faculdades, e 28,3, nos CEFET/IF.
Observa -se que existe diferença estatisticamente significativa, ao nível de 95 \%, no Componente de Conhecimento Específico entre as notas das Faculdades em relação às notas das Universidades e dos CEFET/IF.
Quanto a Categoria Administrativa, observa -se um comportamento semelhante àquele da parte de Formação Geral e à prova, como um todo, ou seja, existe diferença estatisticamente significativa entre as Médias das IES Públicas ( 43,9 ) e as das IES Privadas ( 35,4 ).
Neste caso, também, a maior Média foi obtida por estudantes de IES Públicas de ensino. 6.2 ANÁLISE DAS QUESTÕES OBJETIVAS Esta seção apresenta estatísticas selecionadas e histogramas das Questões Objetivas de Formação Geral ( 6.2.1 ) e de Conhecimento Específico ( 6.2.2 ).
São, também, apresentadas e comparadas as médias das subpopulações caracterizadas por Grande Região. 6.2.1 Componente de Formação Geral Na Tabela 6.7, são apresentadas as Estatísticas Básicas relativas às oito questões objetivas do componente da prova que abrange a Formação Geral dos estudantes.
A Média do Brasil foi 55,9.
A menor Média foi encontrada na região Norte ( 51,2 ), e a maior, na região Sudeste ( 60,7 ).
As demais Médias foram : 55,8, na região Nordeste, 55,6, na região Sul e 57,6, na região Centro-Oeste.
O Desvio padrão do Brasil foi 21,7, sendo o maior Desvio padrão encontrado na região Norte ( 22,1 ), e o menor, na região Sul ( 20,4 ).
Os demais desvios foram : 21,9, na região Nordeste, 21,2, na região Sudeste e 21,6, na região Centro-Oeste.
As Medianas do Brasil, como um todo, e das cinco regiões foram iguais a 57,1.
A nota Máxima 100,0 foi alcançada nas cinco Grandes Regiões.
As notas mínimas ( 0,0 ) foram iguais para todas as regiões, exceto para a região Nordeste ( 14,3 ).
Na Tabela 6.8, são apresentados o Índice de Facilidade e o Índice de Discriminação ( Ponto-Bisserial ) para \textbf{cada} uma das questões objetivas do Componente de Formação Geral.
Quanto ao Índice de Facilidade, foram usadas as seguintes cores para diferenciar o nível de dificuldade da questão :  Azul para as questões classificadas com índice Muito fácil ( \&gt ; =0,86 ), verde para as questões classificadas com índice Fácil ( 0,61 a 0,85 ), amarelo para as questões classificadas com Médio ( 0,41 a 0,60 ), vermelho para as questões classificadas com Difícil ( 0,16 a 0,40 ) e roxo para as questões classificadas com Muito difícil ( \&lt ; =0,15 ).
Já quanto ao Índice de Discriminação, foram usadas as seguintes cores para qualificar a questão :  As questões classificadas com índice Fraco receberam a cor vermelha ( \&lt ; =0,19 ), as classificadas com Médio receberam a cor amarela ( 0,20 a 0,29 ), as classificadas com Bom receberam a cor verde ( 0,30 a 0,39 ) e as classificadas com Muito bom ( \&gt ; =0,40 ) receberam a cor azul.
As questões objetivas do Componente de Formação Geral, segundo o Índice de Facilidade, foram assim avaliadas : das sete questões ( uma foi anulada pela CAA – Comissão Assessora da Área ), nenhuma teve o Índice de Facilidade classificado como Muito fácil, e três questões foram tidas como Fácil situando -se no intervalo entre 0,63 e 0,70 do Índice de Facilidade, ou seja, houve entre 63,0 \% e 70,0 \% de acertos.
Outras três questões foram consideradas com índice de dificuldade Médio, situando -se no intervalo entre 0,51 e 0,59 do Índice de Facilidade, ou seja, houve entre 51,0 \% e 59,0 \% de acertos, enquanto uma questão foi classificada na categoria Difícil, obtendo um índice de 0,25.
Nenhuma questão apresentou menos de 15 \% de acertos, razão pela qual não houve questão classificada como Muito difícil.
E a questão 1 foi anulada pela CAA.
O Índice de Facilidade variou de 0,25 a 0,70 e o de Discriminação, de 0,37 a 0,55.
Seis questões, quanto ao Índice de Discriminação, foram tidas como Muito bom.
E apenas uma tida como Bom.
Na Tabela 6.9, é apresentada a distribuição das questões, levando -se em conta, simultaneamente, a classificação dos dois índices.
As seis questões que tiveram Índice de Discriminação Muito bom figuraram entre três níveis de dificuldade Fácil, Médio e Difícil : três classificadas na categoria Fácil ( as questões 2, 3 e 8 ) do Índice de Facilidade, duas, na categoria Médio ( questões 4 e 6 ), e uma, na categoria Difícil ( questão 5 ).
Em particular, a questão 8 foi a que apresentou o maior poder discriminatório, com índice 0,55, e foi considerada Fácil em termos de facilidade, com uma proporção de 0,63 acertos.
O máximo de acertos foi alcançado pelas questões 2 e 3, com um Índice de Facilidade de 0,70, em ambas.
A questão 5 foi considerada Difícil, com índice de Facilidade de 0,25.
Como dito, anteriormente, a questão 1 foi anulada pela CAA.
No Gráfico 6.4, para exemplificar, analisa -se o comportamento da questão de número 8 de Formação Geral.
Trata -se de uma questão considerada Fácil, em relação à facilidade e a que obteve o segundo maior Índice de Discriminação dessa parte da prova.
Neste gráfico, \textbf{cada} uma das cinco curvas representa o percentual de respostas em determinada alternativa da questão, em função do número de acertos dos estudantes nessa parte da prova ( Formação Geral/Múltipla Escolha ), antes de possíveis eliminações pelo critério Ponto-Bisserial.
A curva em verde corresponde à alternativa B, a alternativa correta para essa questão.
Observa -se que, entre os estudantes com menor número de acertos, nessa parte do exame, a situação mais frequente foi a escolha da alternativa C ( em preto ), incorreta.
Por exemplo, entre os estudantes que acertaram duas questões, 14,3 \% escolheram a alternativa E ( em vermelho ), 17,0 \% escolheram a alternativa D ( em laranja ), 22,4 \% escolheram a alternativa B ( em verde, gabarito ), 11,7 \% escolheram a alternativa A ( em azul ), e 34,1 \%, a C ( em preto ).
Dentre os que acertaram duas respostas entre as questões de múltipla escolha de Formação Geral, 0,4 \% deixou a questão em branco e 0,0 \% marcou mais de uma alternativa, invalidando a questão.
À medida que o número de acertos aumenta, indicando desempenho melhor nessa parte da prova, aumenta, concomitantemente, a proporção de estudantes que selecionaram a alternativa correta B, atingindo 100 \% para os estudantes com sete acertos.
Essa análise permite verificar como a questão discriminou os grupos de desempenho, justificando -se o alto índice ( 0,55 ) obtido na questão.
Cumpre notar que não é possível inferir deste gráfico nem o Índice de Facilidade, que seria uma média da proporção ponderada pela quantidade de estudantes com \textbf{cada} uma das notas, nem o Índice de Discriminação Ponto-Bisserial, por razão equivalente.
No caso extremo, no qual a grande concentração dos acertos dos estudantes fosse abaixo de quatro, o Índice de Facilidade seria obrigatoriamente abaixo de 70,0 \% ( neste exemplo ).
Caso a concentração fosse em seis acertos ou mais, o índice seria obrigatoriamente acima de 90,0 \%.
Os gráficos relativos às demais questões de Formação Geral constam do Anexo I. 6.2.2 Componente de Conhecimento Específico Na Tabela 6.10, são apresentadas as Estatísticas Básicas em relação às questões objetivas do Componente de Conhecimento Específico da prova de Engenharia Florestal por Grande Região.
Como será discutido ainda nesta seção, três questões objetivas do Componente de Conhecimentos Específicos não foram usadas no cômputo das notas por terem sido descartadas pelo critério do Ponto-Bisserial.
Assim, as notas foram calculadas com base em 24 das 27 questões objetivas válidas de Conhecimentos Específicos.
A Média do Brasil deste componente foi 44,2.
A menor Média foi observada na região Norte ( 38,6 ), e a maior, na região Sudeste ( 48,9 ).
O Desvio padrão de todo o Brasil foi 14,0, sendo o menor Desvio padrão encontrado na região Norte ( 12,7 ), e o maior, na região Sudeste ( 14,5 ).
A Mediana de todo o Brasil foi 41,7, o mesmo valor da Mediana encontrada na região Nordeste.
Nas demais regiões as Medianas foram : 37,5, na região Norte, 50,0, na região Sudeste, e 45,8, nas regiões Sul e Centro-Oeste.
A nota Máxima da prova foi 91,7, obtida nas questões objetivas do Componente de Conhecimento Específico, por, pelo menos, um estudante da região Sudeste.
Nas demais regiões, as notas Máximas foram : 75,0, na região Norte, 83,3, nas regiões Nordeste e Sul, e 75,0 na região Centro-Oeste.
A nota Mínima foi zero, obtida por, pelo menos, um estudante da região Sudeste.
Nas demais regiões, as notas Mínimas foram : 4,2, nas regiões Norte e Nordeste ; 8,3, na região Sul ; e 12,5, na região Centro-Oeste.
Na Tabela 6.11, são apresentados os Índices de Facilidade e de Discriminação ( Ponto-Bisserial ) das questões objetivas do Componente de Conhecimento Específico, para os estudantes de Engenharia Florestal.
Para facilitar a diferenciação das questões, serão usadas as mesmas cores da Tabela 6.8 para as diferentes classificações dos Índices de Facilidade e de Discriminação.
Observa -se que as análises dizem respeito às 27 questões válidas.
A partir do Índice de Facilidade obtido, pode -se concluir que mais da metade das questões objetivas da prova foram consideradas, pelo menos, Difícil : das 27 questões válidas, 14 foram classificadas como Difícil ou como Muito difícil.
Três questões foram classificadas como Fácil, e outras dez consideradas como Médio.
Nenhuma questão foi classificada como Muito fácil.
Já quanto aos Índices de Discriminação das questões objetivas do Componente de Conhecimento Específico da prova, obtém -se como resultado a seguinte classificação : 11 das 27 questões foram consideradas como Bom, enquanto nenhuma delas teve Índice de Discriminação Muito bom.
Dentre as demais, 13 delas foram classificadas como Médio, e outras três, como Fraco, sendo 16, por conseguinte, a quantidade de questões nos dois patamares mais baixos de discriminação.
Constata -se, assim, que a prova – no que se refere ao Componente de Conhecimento Específico – possuía média capacidade de discriminar entre aqueles que dominam ou não o conteúdo.
O Índice de Facilidade variou de 0,07 a 0,83 e o de Discriminação, de 0,03 a 0,39.
Na Tabela 6.12, é apresentada a distribuição das questões levando -se em conta, simultaneamente, a classificação dos dois índices.
Duas questões alcançaram os maiores Índices de Discriminação, classificadas com índice Bom, as de número 13 e 17.
Seus índices foram 0,60 e 0,48, respectivamente, e quanto ao Índice de Facilidade, ambas foram classificadas como Médio, com 60 \% e 48 \%, respectivamente, dos estudantes marcando a opção correta.
A questão de número 28 foi a mais difícil entre as 27 questões específicas válidas, com baixo Índice de Facilidade, apenas 7 \% de acertos.
Essa questão apresentou poder discriminatório muito baixo, 0,03, o que comprova ter sido a mais difícil para os estudantes.
Destacam -se, também, as questões 27 e 35, com Índice de Facilidade 0,28 e 0,18, respectivamente, o que, em termos percentuais, correspondem a 28 \% e 18 \% de estudantes que responderam acertadamente, respectivamente.
Já 0,19 e 0,05 foram os seus Índices de Discriminação.
Tais questões foram, portanto, pelo critério Ponto-Bisserial, consideradas inadequadas.
Por isso, as questões 27, 28 e 35 foram eliminadas do cômputo da nota final.
A título de exemplo das análises do comportamento das questões objetivas, no Gráfico 6.5, analisa -se a questão 13 do Componente de Conhecimento Específico.
Essa questão foi considerada, pelas respostas dos estudantes avaliados na prova, como uma questão com nível de facilidade Médio, 0,60, ou seja, 60 \% dos estudantes assinalaram, acertadamente, a opção A, correspondente ao gabarito.
Como já comentado, seu Índice de Discriminação foi igual a 0,39, classificado como Bom, o maior valor de discriminação.
Neste gráfico, \textbf{cada} uma das cinco curvas representa o percentual de respostas em determinada alternativa da questão 13, em função do número de acertos dos estudantes nessa parte da prova, antes de possíveis eliminações de questões pelo critério Ponto-Bisserial.
A alternativa correta A, representada no gráfico pela curva em azul, foi escolhida em maiores proporções pelos estudantes com desempenho melhor nessa parte da prova.
Já as alternativas incorretas, também denominadas distratores, foram selecionadas, principalmente, por aqueles com notas mais baixas.
No caso de Engenharia Florestal, como se observa no eixo horizontal do Gráfico 6.5, nenhum estudante acertou mais do que 22 questões dentre as 27 questões objetivas, considerando -se o gabarito original.
Assim, mesmo com a eliminação das três questões com Índice de Discriminação Fraco do cômputo da nota, nenhum participante atingiu a nota Máxima, igual a 100,0, nessa parte da prova ( visto na Tabela 6.10 ).
Observa -se que a soma não é 100 \%, por causa das questões não respondidas ou com mais de uma opção marcada.
Aqueles com nota zero, na sua quase totalidade, deixaram esta questão em branco ou marcaram mais de uma alternativa, comportamento considerado inválido.
A proporção de estudantes que selecionaram a resposta correta A aumenta gradativamente, chegando a atingir 100 \% para 20 acertos ou mais, enquanto a proporção dos que escolheram alternativas incorretas decai, a partir de três ou quatro acertos, em função do número de acertos nessa parte da prova.
Os gráficos relativos às demais questões do Conhecimento Específico constam do Anexo I. 6.3 ANÁLISE DAS QUESTÕES DISCURSIVAS Esta seção apresenta estatísticas selecionadas e histogramas das Questões Discursivas de Formação Geral ( 6.3.1 ) e Conhecimento Específico ( 6.3.2 ).
São, também, apresentadas e comparadas as médias de subpopulações, caracterizadas por Grande Região.
Cumpre notar que nem todas as questões passam pelo crivo dos corretores.
As respostas depois de digitalizadas passam por um pré-processamento de identificação automático.
Algumas destas questões são definidas como um “ branco automático ” ( que não são enviadas para os professores para correção, a não ser que na segunda fase, a de inspeção visual, seja constatada uma resposta ).
O critério usado neste pré-processamento, é verificar para \textbf{cada} questão, a partir da quantidade de PIXELS, o que provavelmente, é uma resposta em branco.
A partir da separação destes documentos que foram considerados em branco, eles são enviados para uma verificação visual, por uma equipe de colaboradores, devidamente treinados.
Deste modo, quaisquer erros, deste pré-processamento, são encontrados, e as questões preenchidas, são enviadas para a correção dos professores.
Eventualmente, algumas respostas em branco, não são detectadas neste pré-processamento automático.
A presença de pequenos riscos, ou manchas, podem fazer o programa não identificar uma questão como um possível branco.
Estas são enviadas para a correção, mas são classificadas como “ branco ”, pelos professores quando da correção.
No caso da Área de Engenharia Florestal, foram identificados como brancos “ automáticos ”, 277 respostas de questões discursivas de Conhecimento Específico e 221 de Formação Geral.
Pelos professores corretores foram identificadas mais 145 e 27 questões em branco, respectivamente, para Conhecimento Específico e para Formação Geral.
Uma outra situação a ser mencionada é a ocorrências de atendimento a portadores de necessidades especiais.
O anexo X apresenta, para o exame como um todo e para a área de Engenharia Florestal, por tipo de deficiência, o protocolo usado para permitir a correção das respostas dos alunos em \textbf{cada} situação. 6.3.1 Componente de Formação Geral As análises dos resultados de desempenho dos estudantes de Engenharia Florestal, nas duas questões discursivas relativas a Formação Geral, encontram -se na Tabela 6.13 e no Gráfico 6.6.
Aqui, analisa -se a média das notas de \textbf{cada} questão discursiva de Formação Geral que leva em conta as correções realizadas por duas bancas, uma que analisa e pontua aspectos relativos ao conteúdo das respostas, seguindo o padrão de respostas divulgado pelo Inep ( ver Anexo VIII ), outra que avalia o desempenho linguístico expresso pela redação das respostas, seguindo um padrão de respostas específico da área de Língua Portuguesa ( ver Anexo VIII ).
A nota de \textbf{cada} questão é composta de 80 \% da nota da correção de conteúdo e 20 \% da nota de desempenho linguístico.
Na Tabela 6.13, observa -se que a nota Média nesse conjunto de questões foi inferior à obtida nas objetivas.
Os estudantes de todo o Brasil obtiveram, em Formação Geral, Média 55,9, nas questões objetivas e 34,0, nas questões discursivas.
Pode -se notar, também, que o Desvio padrão nesse conjunto de questões foi menor do que o obtido nas objetivas : 21,7, nas questões objetivas e 19,6, nas questões discursivas.
A maior Média foi obtida na região Sudeste ( 39,6 ), e a menor, na região Norte ( 29,4 ).
A Mediana de todo o Brasil, neste componente, foi 32,5.
Nas regiões Norte e Centro-Oeste, a Mediana foi menor ( 28,5 e 31,5 respectivamente ), e, nas demais regiões foi igual ou maior que o valor para o Brasil : Nordeste ( 32,8 ), Sudeste ( 39,0 ), Sul ( 32,5 ).
A nota Máxima ( 95,0 ) foi obtida na região Centro-Oeste, sendo 93,5, a Máxima na região Norte ; 85,5, na Nordeste, 91,5, na Sudeste e 84,5, na Sul.
A nota Mínima ( 0,0 ) foi a mesma em todas as regiões do Brasil.
No Gráfico 6.6, \textbf{está} representada a distribuição das notas nas questões discursivas do Componente de Formação Geral.
A moda dessa distribuição ocorre no intervalo ( 30 ; 40 ] com distribuição de 18,9 \% do total de notas.
Destacam -se, também, os intervalos ( 20 ; 30 ] e ( 40 ; 50 ] com distribuições de 16,9 \% e 16,7 \% do total de notas, respectivamente.
Na sequência, os resultados verificados para \textbf{cada} uma das questões discursivas de Formação Geral serão apresentados, estabelecendo -se relações com os conteúdos abordados em \textbf{cada} uma delas.
Os comentários das Bancas de docentes corretores, a respeito do observado na correção das respostas dos estudantes, suas impressões e conclusões serão apresentados junto à análise de \textbf{cada} questão.
Cumpre esclarecer que, tendo em vista que as questões discursivas de Formação Geral são padronizadas, ou seja, constam de todas as provas, os comentários da Banca são os mesmos para todas as carreiras acadêmicas, sendo direcionados a todos os estudantes que participaram do Enade/2019.
A seguir, será analisado o desempenho linguístico dos estudantes da Área de Engenharia Florestal nas duas questões discursivas de Formação Geral do Enade/2019, apresentando os resultados obtidos e os comentários da banca de correção para \textbf{cada} questão. 6.3.1.1 Análise de Conteúdo da Questão Discursiva 1 do Componente de Formação Geral Os dados de Engenharia Florestal, obtidos a partir das respostas à questão 1, encontram -se na Tabela 6.14 e no Gráfico 6.7.
Nessa questão – de desempenho abaixo do da outra questão de Formação Geral – os estudantes de todo o Brasil, tiveram Média 26,7.
A maior Média para a questão 1 foi obtida na região Sudeste ( 34,0 ), e a menor, na região Norte ( 22,0 ).
Quanto à variabilidade das notas, o Desvio padrão de todo o Brasil foi 31,0.
O menor Desvio padrão foi obtido na região Norte ( 28,4 ), e o maior Desvio padrão foi obtido pela região Nordeste ( 33,9 ).
A Mediana do Brasil, como um todo, foi 20,0, sendo igual somente nas regiões Nordeste e Centro-Oeste.
Nas demais regiões, as medianas foram zero, na Norte e na Sul, e 40,0, na Sudeste.
As notas Máxima e Mínima da questão discursiva 1 foram, sem exceção, as mesmas para todas as regiões do Brasil, respectivamente, 100,0 e zero.
No Gráfico 6.7, mostra -se a distribuição das notas na questão discursiva 1 do Componente de Formação Geral.
Observa -se que 5,8 \% dos participantes deixaram em branco a resposta a essa questão.
Somados à frequência dos que, tendo respondido, receberam nota zero, chega a 46,1 \% o contingente daqueles que zeraram questão.
Sem considerar a nota zero, o intervalo modal seria ( 40 ; 50 ], representado 14,7 \% das notas alocadas nesse intervalo, seguido de perto do intervalo ( 20 ; 30 ]. receberam nota zero, chega a 46,1 \% o contingente daqueles que zeraram questão.
Sem considerar a nota zero, o intervalo modal seria ( 40 ; 50 ], representado 14,7 \% das notas alocadas nesse intervalo, seguido de perto do intervalo ( 20 ; 30 ]. 6.3.1.2 Comentários sobre a correção de Conteúdo das respostas à Questão Discursiva 1 No enunciado da Questão 1, que era constituída por dois textos e um comando, abordavam -se dois itens da lista de temas de referência estabelecidos pelo Inep, pela Portaria nº 510, de 31 de maio de 2019, que dispõe sobre o componente de Formação Geral : ( VIII ) meio ambiente, sustentabilidade e intervenção humana e ( IX ) cidades, habitação e qualidade de vida.
No primeiro texto, apresentavam -se estatísticas sobre desastres naturais ocorridos no Brasil e a eles associavam -se eventos recorrentes, tais como estiagens, secas, inundações bruscas e alagamentos.
O segundo texto complementava o primeiro, nele constava um relatório da ONU no qual é identificada a necessidade de se reduzirem os níveis existentes de riscos que favorecem os desastres, apresentando o fortalecimento da resiliência social, ambiental e econômica como uma das soluções para que as cidades consigam conviver com esses fenômenos naturais.
No comando, por outro lado, não se solicitava reflexão sobre os riscos de futuros desastres, temática do segundo texto.
Solicitava -se reflexão sobre ações após a ocorrência de um desastre : “ duas propostas de intervenção no âmbito da sustentabilidade socioambiental de modo a contemplar ações de recuperação ou de restauração após a ocorrência de desastres ”.
Assim, apesar de os textos de contextualização serem claros e conterem informações relevantes – fatos de conhecimento público e de grande divulgação –, havia uma diferença de abordagem entre os dois e o comando em relação à linha do tempo da ocorrência de desastres.
Esse fato aumentou o grau de dificuldade da questão.
Entretanto, os conhecimentos exigidos e o nível de profundidade esperado pelo padrão de resposta eram compatíveis com a formação universitária.
Ao se exigirem propostas de intervenção envolvendo a sustentabilidade socioambiental, a questão possibilitou a avaliação, em primeiro lugar, da capacidade de elaborar e apresentar propostas de ação e de intervenção, a partir da análise de um contexto, e, em segundo lugar, da capacidade de propor soluções viáveis e inovadoras na resolução de uma situação-problema.
Na solicitação do comando \textbf{estava} implícita a exigência de redação de um texto em que se apresentasse \textbf{cada} proposta, o que deveria incluir informações suficientes para sua compreensão, ou seja, como as ações seriam concretizadas, por quem seriam realizadas ou que benefícios trariam à população ou à área atingida.
Na questão, exigia -se, portanto, a capacidade de o estudante expor e argumentar, de modo a defender ou justificar suas propostas, de forma coerente e articulada.
Além disso, pelo conteúdo da resposta era possível avaliar o estágio de comprometimento do estudante com as questões sociais e ambientais.
A dicotomia entre o antes e o após desastre, criada no confronto entre os textos e o comando da questão, permitiu aquilatar também a aptidão do estudante para ler e interpretar textos não lineares, em que são exigidas atenção e percepção de diferenças sutis.
No entanto, tal dicotomia gerou respostas que não \textbf{estavam} em conformidade com as linhas de ação apontadas no padrão de resposta, respostas erradas por terem foco em ações de prevenção.
O padrão de respostas foi apresentado em linguagem simples, compatível com o nível dos formandos.
O padrão tinha flexibilidade e extensão, visto que apresentava inúmeras ações de recuperação e de restauração, no âmbito da sustentabilidade socioambiental, sem exauri -las, como passíveis de abordagem pelos estudantes em suas respostas.
Também não havia exigência de conhecimentos específicos sobre qualquer matéria, ficando restrito aos limites de conceitos básicos sobre meio ambiente, sustentabilidade e intervenção humana em econômica como uma das soluções para que as cidades consigam conviver com esses fenômenos naturais.
No comando, por outro lado, não se solicitava reflexão sobre os riscos de futuros desastres, temática do segundo texto.
Solicitava -se reflexão sobre ações após a ocorrência de um desastre : “ duas propostas de intervenção no âmbito da sustentabilidade socioambiental de modo a contemplar ações de recuperação ou de restauração após a ocorrência de desastres ”.
Assim, apesar de os textos de contextualização serem claros e conterem informações relevantes – fatos de conhecimento público e de grande divulgação –, havia uma diferença de abordagem entre os dois e o comando em relação à linha do tempo da ocorrência de desastres.
Esse fato aumentou o grau de dificuldade da questão.
Entretanto, os conhecimentos exigidos e o nível de profundidade esperado pelo padrão de resposta eram compatíveis com a formação universitária.
Ao se exigirem propostas de intervenção envolvendo a sustentabilidade socioambiental, a questão possibilitou a avaliação, em primeiro lugar, da capacidade de elaborar e apresentar propostas de ação e de intervenção, a partir da análise de um contexto, e, em segundo lugar, da capacidade de propor soluções viáveis e inovadoras na resolução de uma situação-problema.
Na solicitação do comando \textbf{estava} implícita a exigência de redação de um texto em que se apresentasse \textbf{cada} proposta, o que deveria incluir informações suficientes para sua compreensão, ou seja, como as ações seriam concretizadas, por quem seriam realizadas ou que benefícios trariam à população ou à área atingida.
Na questão, exigia -se, portanto, a capacidade de o estudante expor e argumentar, de modo a defender ou justificar suas propostas, de forma coerente e articulada.
Além disso, pelo conteúdo da resposta era possível avaliar o estágio de comprometimento do estudante com as questões sociais e ambientais.
A dicotomia entre o antes e o após desastre, criada no confronto entre os textos e o comando da questão, permitiu aquilatar também a aptidão do estudante para ler e interpretar textos não lineares, em que são exigidas atenção e percepção de diferenças sutis.
No entanto, tal dicotomia gerou respostas que não \textbf{estavam} em conformidade com as linhas de ação apontadas no padrão de resposta, respostas erradas por terem foco em ações de prevenção.
O padrão de respostas foi apresentado em linguagem simples, compatível com o nível dos formandos.
O padrão tinha flexibilidade e extensão, visto que apresentava inúmeras ações de recuperação e de restauração, no âmbito da sustentabilidade socioambiental, sem exauri -las, como passíveis de abordagem pelos estudantes em suas respostas.
Também não havia exigência de conhecimentos específicos sobre qualquer matéria, ficando restrito aos limites de conceitos básicos sobre meio ambiente, sustentabilidade e intervenção humana em emergências provocadas por desastres, além de atualização com assuntos amplamente tratados na mídia nacional.
O padrão cobriu diversas áreas de atividades e categorizou -as em ações psicossociais, econômicas e sociais, ambientais, de infraestrutura e sistêmicas, com destaque para as ocorrências de estiagens e secas, porque são fenômenos praticamente permanentes em diversas regiões do país.
Esse universo de intervenções criou amplo espaço para o enquadramento das respostas dadas pelos estudantes, incluindo enfoques variados do ponto de vista das diferentes formações e especificidades das áreas dos participantes do Enade/2019.
O conteúdo das respostas analisadas na correção foi bem coberto pelas propostas sugeridas no padrão.
Da mesma forma, as respostas corretas cobriram praticamente todo o espectro levantado pelo padrão.
Na maior parte das respostas que não atenderam ao padrão, propunham -se ações de prevenção, ou seja, procedimentos anteriores aos desastres e, portanto, equivocadas.
Nas respostas alinhadas com o padrão, preferencialmente, havia sugestões de ações, dentre as quais se destacaram : recuperação da flora nativa, controle da qualidade da água, organização e treinamento de grupos de voluntários para ações de prevenção de surtos e epidemias, orientação para evitar novos desabamentos e redução de efeitos dos alagamentos, concessão de benefícios como aluguel social e financiamento para recuperação de residências ou reativação de negócios, projetos de realocação de desabrigados e suporte psicológico, médico e social às pessoas e comunidades afetadas pelas calamidades.
Nas respostas envolvendo situações de estiagem e seca, apresentavam -se ações que tanto poderiam ser consideradas de prevenção quanto de recuperação, pelas características sistêmicas desses desastres.
Por isso, propostas como distribuição de água por caminhões pipa ou perfuração de poços foram consideradas corretas.
Apesar de fora do contexto da questão, foram bastante citadas sugestões de penalização dos supostos culpados pelo desastre com multas, bem como recompensa às pessoas ou empresas que prestassem socorro nessas situações.
Por outro lado, apesar de presente no padrão, ações de resgate do patrimônio histórico e cultural praticamente não foram citadas como resposta.
Em um conjunto grande de respostas, os estudantes se limitaram a apresentar apenas uma proposta e, nesse caso, preferencialmente, de ajuda a desabrigados ou de recuperação de matas ou de ações vinculadas à seca e à estiagem.
Dentre as respostas consideradas corretas, muitas eram pobres na argumentação e, por vezes, confusas, apesar de \textbf{estarem} em conformidade com o padrão.
Foi percebida \textbf{certa} dificuldade para se organizar a argumentação no contexto do comando.
Em muitas respostas, citavam -se ações isoladas sem que se caracterizasse o seu desdobramento.
Como já comentado, grande parte dos estudantes, em desacordo com o comando da questão, fixou -se em ações de prevenção, anteriores aos desastres.
Houve respostas de excelente conteúdo, mas totalmente enquadradas no contexto preventivo.
Esse viés de prevenção também incluiu respostas fracas em que se parafraseavam os textos do enunciado, mencionando -se de forma superficial inundação, alagamentos, seca e redução dos níveis de risco.
A interpretação incorreta pode ter sido induzida pelo enunciado, mas o principal indício parece ser de falta de atenção na leitura do comando.
As maiores fragilidades na formação dos estudantes, apontadas pelos corretores, foram a leitura e interpretação da questão e a capacidade de expressão linguística.
As respostas com geração de ideias próprias ocorreram em escala bem menor do que aquelas em que os estudantes usaram do recurso da paráfrase pobre sobre o texto do enunciado, revelando dificuldade para criar e articular ideias a respeito dos temas da questão.
Além disso, pelas sugestões que envolviam aspectos econômicos e financeiros revelou -se um nível razoável de desconhecimento das atribuições e responsabilidades das diversas instâncias do governo e de como são encaminhados os recursos necessários ao socorro de populações atingidas por desastres.
Como aspectos relevantes ainda não considerados, podem ser citados o baixo percentual de protestos e a falta de viés crítico nas respostas, talvez pelo fato de o comando ser bastante objetivo. 6.3.1.3 Análise de Conteúdo da Questão Discursiva 2 do Componente de Formação Geral A Tabela 6.15 mostra que o desempenho médio dos estudantes, na questão discursiva 2 ( Média 31,0 ), foi acima do obtido na questão discursiva 1 ( Média 26,7 ).
A região Sudeste foi aquela cuja Média, também nessa questão, foi maior ( 35,9 ), assim como a de menor Média foi a da região Norte ( 26,1 ).
Quanto à variabilidade das notas, o Desvio padrão de todo o Brasil foi 29,4, um pouco inferior ao obtido na questão discursiva 1 ( 31,0 ).
O maior desvio nessa questão foi obtido na região Sudeste ( 30,4 ), enquanto o menor foi obtido na região Nordeste ( 27,8 ).
A Mediana de todo o Brasil foi 25,0, a mesma de quatro regiões.
A exceção foi a região Sudeste, que obteve Mediana 50,0.
As notas Máxima ( 100,0 ) e Mínima ( 0,0 ) foram as mesmas em todas as regiões, sem exceção.
No Gráfico 6.8, mostra -se a distribuição das notas na questão discursiva 2 do Componente de Formação Geral.
Observa -se que, 7,6 \% dos estudantes deixaram a questão em branco.
Chega a 38,8 \%, se agregado à frequência daqueles que, tendo respondido, receberam nota zero.
Sem considerar a nota zero, o intervalo de maior frequência foi ( 50 ; 60 ], seguido do intervalo ( 70 ; 80 ].
Nota -se, ainda, que as notas ficaram levemente menos dispersas em comparação às da questão discursiva 1, o que pode ser constatado, também, pela comparação do desvio padrão das notas da questão discursiva 2 ( 29,4 ) e o das notas da questão discursiva 1 ( 31,0 ). 6.3.1.4 Comentários sobre a correção de Conteúdo das respostas à Questão Discursiva 2 Na Questão 2, o tema era a identificação de impactos científicos e econômicos do desenvolvimento do Brasil na área de Ciência e Tecnologia ( C\&T ), bem como o papel dos principais atores responsáveis por impulsionar esse desenvolvimento.
O tema escolhido era relevante por contemplar os referenciais divulgados pelo Inep para esse exame : Ciência, tecnologia e inovação – Estado, sociedade e trabalho, assuntos supostamente acessíveis aos formandos do nível universitário brasileiro de qualquer área.
A análise da questão demonstra sua conformidade com as diretrizes de avaliação de competências e habilidades dos estudantes.
A explicitação de ganhos vinculados a dois campos diferentes, no contexto de C\&T, e como consequência de colaboração interinstitucional, proporciona a verificação da capacidade não só de ler e interpretar textos, como de narrar, de forma consistente e articulada, o resultado de reflexão sobre necessidades em contextos diversos.
Além disso, as descrições de articulações e escolhas de ganhos que se julgassem relevantes permitiram que se avaliasse o nível de maturidade do conhecimento sobre os sistemas científicos, tecnológicos e econômicos brasileiros.
O enunciado era claro, e sua linguagem, compatível com o que se espera de formandos de cursos superiores.
Era constituído por uma figura e um texto de contextualização que auxiliavam na compreensão do que foi solicitado pelo comando.
No texto de apresentação, comparava -se a posição do Brasil com a de outros países em indicadores de produção e utilização de conhecimento e de novas tecnologias.
Na figura ilustrava -se a relação entre as principais entidades presentes no contexto de C\&T – as instituições de ensino superior, as empresas e o setor público.
Havia o pressuposto, portanto, de que os egressos fossem capazes de sustentar ideias articuladas e níveis de argumentação, a partir dos elementos textuais fornecidos pela questão, a saber, a figura e o texto, elementos de amparo à elaboração da resposta.
No comando, interligando -se texto e figura, pedia -se para o estudante citar dois ganhos possíveis para o campo científico e dois ganhos possíveis para o campo econômico, resultantes de uma boa articulação entre os entes representados na figura.
Aparentemente, a tarefa solicitada seria simples.
Entretanto, exigia reflexão do estudante em três aspectos : ( i ) a exigência de um processo de articulação entre, pelo menos, duas das entidades explicitadas na figura ; ( ii ) a criação de duas propostas de ganhos em dois setores diferentes, o científico e o econômico, ambos no contexto de C\&T ; e ( iii ) que os ganhos fossem consequência da colaboração de, pelo menos, dois dos entes.
Essa singularidade do enunciado, elevou a dificuldade da questão.
O comando permitia duas possibilidades de interpretação da articulação entre os entes para a estruturação das respostas.
Uma seria o estudante interpretar que seria necessário explicitar em sua resposta o inter-relacionamento gerador dos ganhos ; a outra, induzida pela figura, seria considerar que a articulação entre entidades já \textbf{estava} dada pelo enunciado, deixando ao estudante apenas o registro dos ganhos.
É presumível que a intenção inserida no comando, e refletida no padrão de resposta, fosse a primeira opção, mas a segunda interpretação era plausível.
O padrão corroborava a primeira opção, pois parecia esperar dos formandos a identificação de uma relação de causa-efeito da ação conjunta de dois ou mais entes e do relato das consequências científicas e econômicas dessa ação conjunta.
O padrão de resposta era abrangente e redigido em linguagem simples, apropriada, com nível de profundidade coerente com a formação geral esperada dos respondentes.
Foi apresentado na forma de duas listas extensas, com três temas científicos e oito econômicos, expressos de forma genérica, \textbf{cada} tema completado com exemplos de ganhos passíveis de citação pelos estudantes.
A variedade de temas abordados no padrão proporcionou a flexibilidade e a extensão necessárias para orientar critérios de avaliação.
As respostas mantiveram -se, em sua grande maioria, no escopo antecipado pelo padrão de respostas e no universo de linguagem adequado ao cenário de C\&T da questão.
Nesse contexto da linguagem, foi peculiar constatar a existência de grupos de respostas semelhantes e reveladoras da área de origem dos estudantes.
Por exemplo, desenvolvimento de vacinas, pele de tilápia no tratamento de queimaduras e criação de projetos de máquinas industriais de maior produtividade.
A redação, entretanto, foi quase sempre mais simples do que o esperado pelo padrão de resposta, sem desenvolvimento de argumentação mais aprofundada e, às vezes, sem qualquer tipo de argumentação.
Há dois motivos para esse discurso mais direto dos estudantes.
O mais importante veio do próprio comando da questão, em que constava o verbo citar ; liberando a resposta de um compromisso maior com a dissertação.
O segundo, talvez induzido pela figura, veio da compreensão de que a articulação já \textbf{estaria} pressuposta e que bastava a apresentação dos ganhos.
As respostas mais bem avaliadas mantiveram -se no contexto do padrão, mas seguiram duas vertentes ; nelas, os estudantes consideravam ou não a articulação como pré-estabelecida.
Aquelas em que os estudantes fizeram a associação no próprio texto da resposta tinham um conteúdo mais rico e mais bem desenvolvido.
Outra característica das respostas corretas, na conformidade com o padrão de resposta, foi a separação precisa entre os ganhos científicos e econômicos, mantendo -se o contexto de C\&T do texto de apresentação.
Nas respostas equivocadas, seguindo essas duas linhas de interpretação, os estudantes não fugiram dos temas existentes no padrão de resposta, mas não os exploraram adequadamente.
Foram vários os motivos.
Grande parte formou um conjunto de respostas muito sucintas, que não chegavam a qualificar, com clareza, os ganhos como científicos ou econômicos.
Outro conjunto misturou conceitos de educação com pesquisa, ciência com produção, economia com emprego ou trabalho, dificultando a caracterização do ganho e afastando a resposta do padrão.
As propostas de ganhos com viés econômico geraram maior quantidade de não conformidades do que aquelas que trataram do viés científico.
Isso foi resultado de uma dificuldade demostrada pelos estudantes em evidenciar o vínculo do ganho econômico com o contexto de C\&T, quase sempre apresentado com exemplos de vantagem obtida ou concedida pelo setor público a representantes das outras entidades apontadas no enunciado, tais como : aumento de emprego, mais vagas nas universidades ou redução de impostos.
A análise das respostas revelou que, na maioria delas, os ganhos eram descritos de forma genérica, sem qualquer articulação ou exposição de motivos, sem indicar como ou de onde surgiam, sem especificar o contexto do assunto ou da área de conhecimento envolvida.
Na maioria dos casos, a ideia de citação levou à não redação de frases, mas à menção de palavras ou expressões.
Exemplos frequentes foram : aumento de pesquisa, melhoria de conteúdo, melhoria de estrutura.
Mesmo relevando a interpretação de que essa colaboração já vinha do enunciado, foi surpreendente que em grande parte dessas respostas não se descrevesse como alguma cooperação seria necessária ou deveria ter sido estabelecida para a geração desses ganhos.
Considerar a articulação vinda do enunciado não dispensava essa revelação.
Por outro lado, houve respostas em que estudantes discutiram a importância da articulação, mas, curiosamente, sem citar os ganhos gerados por elas.
Essa dúvida sobre a articulação ficou bem caracterizada nas respostas em que o ganho descrito não vinha de uma parceria, mas simplesmente da ação de uma das entidades em benefício de outra.
No conjunto das respostas sobre ganhos no campo científico, houve prevalência de citações sobre a necessidade de investimentos em pesquisa.
Essa é uma demanda real, bastante discutida no meio acadêmico.
Assim, respostas que apresentaram uma estrutura do tipo “ maiores investimentos/financiamentos em pesquisas ” e valorização dos pesquisadores foram muito comuns.
Geralmente, elas se complementavam com “ pesquisas realizadas pelas IES ”.
Todavia, a argumentação inconsistente dos estudantes evidenciou que existe dificuldade não só em distinguir as funções de pesquisa das de educação, como também a respeito de como ocorrem os processos científicos e tecnológicos.
Poucos demonstraram compreender que os resultados obtidos na área de C\&T são incorporados pelos diversos campos do saber ao longo do tempo e que se consolidam nas técnicas e procedimentos apresentados nos cursos de graduação.
É compreensível que grande parte dos egressos não tenha se envolvido e não venha a seguir o caminho da ciência ; além disso, muitos deles fizeram cursos em EAD ou cursos de graduação tecnológica, não tendo a chance de conviver com esse contexto de C\&T.
No entanto, pelas próprias diretrizes do exame, supõe -se que a formação universitária deveria contemplar a compreensão dessas atividades.
Em uma visão incompleta da ciência, dois fatos importantes ainda foram observados, ambos equivocados : o entendimento tácito de que o desenvolvimento da pesquisa científica é de responsabilidade única das IES e o vínculo desse desenvolvimento a bolsas em cursos de mestrado e doutorado e até mesmo de iniciação científica.
Nessa linha, em quantidade significativa de respostas exemplificou -se o investimento na educação como ganho científico.
No que tange ao item sobre os aspectos econômicos, em boa parte das respostas, os estudantes não estabeleceram relação com o contexto científico demandado pelos elementos da questão.
Nas respostas, não se associavam os ganhos econômicos a qualquer processo científico ou tecnológico.
Muitas vezes, refraseavam a ideia já apresentada como resposta para o campo científico.
Foi majoritário o raciocínio simplista de que ganho econômico é igual a lucro.
Dessa forma, chegaram a citar a “ adoção de mão de obra de estagiários, por elas serem mais baratas que os empregados formais ” como ganho econômico.
Essa concepção de lucro apareceu bastante na proposta de produção e venda de patentes, inclusive para o exterior, talvez porque a palavra patente tenha sido mencionada no texto do enunciado.
Apresentado como resultado das três combinações possíveis de parceria, muitas vezes, esse ganho com patente era usado, na mesma resposta, como ganho científico.
Curiosamente, a desburocratização da sua concessão também foi apresentada como ganho econômico.
Outro ganho econômico bastante citado, por melhorar a lucratividade das empresas, foi a capacitação profissional realizada nas IES, proposta em que há um beneficiário e um doador, mas não há a parceria exigida pelo comando, nem o vínculo com o campo científico tecnológico.
Além disso, as sugestões que envolviam os ganhos econômicos revelaram desconhecimento das atribuições do governo e de como são encaminhados os recursos necessários ao estímulo às atividades cientificas ou econômicas.
A ausência de citações envolvendo órgãos de fomento foi reveladora desse aspecto.
As respostas dos estudantes e o desempenho global retratado na concentração em notas baixas revelam que a formação de nível superior do Brasil, ao menos nos cursos envolvidos nesta edição do Enade, não tem garantido a compreensão do papel da ciência, do papel dos atores que a promovem e de suas consequências socioeconômicas.
Observou -se um desconhecimento generalizado do processo como um todo.
Essa lacuna prejudica tanto o processo de renovação, como o entendimento de que a pesquisa é atividade que requer investimento alto e contínuo, e cujo tempo de retorno destoa do que é praticado no mercado corporativo.
A falta de compreensão do sistema de C\&T, mesmo em nível geral, levou a um número expressivo de respostas em que se indicaram apenas relações de causa e efeito de nível extremamente superficial, do tipo “ a relação entre A e B impacta positivamente a produção científica aumentando as pesquisas ”.
Tal resposta não demonstra compreensão sobre o papel de \textbf{cada} ente, nem de como a relação entre eles se desenvolve para impactar a produção científica. 6.3.1.5 Considerações Finais sobre a correção de Conteúdo das Questões Discursivas do Componente de Formação Geral O principal aspecto detectado nas duas questões foi a dificuldade de os estudantes interpretarem corretamente o que \textbf{estava} sendo solicitado no comando da questão, causada ou por desatenção ou por desconhecimento mais profundo do contexto.
Esse fato foi o principal responsável pelo resultado aquém do esperado para o nível formativo em termos de pontuação.
Os textos daqueles que responderam em conformidade com o padrão demonstraram um universo de argumentação muito limitado, pobre em vocabulário.
O recurso mais comum foi buscar nos textos introdutórios os elementos para a construção de suas frases, o que criou uma peculiaridade interessante, com grupos de respostas muito parecidas.
Encadear e relacionar argumentos também foi uma deficiência identificada pelos docentes da banca de correção.
Os argumentos, quando apresentados, eram redigidos de forma confusa e sem constituir unidade.
Os estudantes apontaram como viés a preocupação social sobre os problemas abordados ; são bons exemplos as propostas de amparo psicológico e material para as vítimas de desastres e o desenvolvimento de vacinas.
Por outro lado, apontaram deficiências consideráveis na compreensão dos papéis que o setor público, as empresas e as IES desempenham na sociedade e no conhecimento tanto do sistema brasileiro de C\&T como da distinção e da articulação entre ensino, pesquisa e extensão. 6.3.1.6 Análise de Língua Portuguesa da Questão Discursiva 1 do Componente de Formação Geral Os dados de Engenharia Florestal, obtidos a partir das respostas à questão discursiva 1 do Componente de Formação Geral, no que tange à Língua Portuguesa, encontram -se na Tabela 6.16 e no Gráfico 6.9.
Nesse aspecto, os estudantes de todo o Brasil obtiveram Média 55,1.
A maior Média em relação à Língua Portuguesa foi obtida na região Nordeste ( 57,0 ), e a menor, na região Norte ( 52,2 ).
Quanto à variabilidade das notas, o Desvio padrão de todo o Brasil foi 20,2.
O menor Desvio padrão foi obtido na região Sudeste ( 17,7 ) e o maior Desvio padrão foi obtido na região Norte ( 22,2 ).
A Mediana das notas de Língua Portuguesa na questão 1 foi 60,0 para o Brasil, como um todo, a mesma obtida nas regiões Nordeste, Sudeste e Sul.
A Mediana das regiões Norte e Centro-Oeste foi 55,0.
A nota Máxima para todo o Brasil foi 90,0, com, pelo menos, um estudante tirando essa nota em todas as regiões, sem exceção.
Já a nota Mínima foi zero em todas as regiões do país.
No Gráfico 6.9, mostra -se a distribuição das notas de Língua Portuguesa do Componente de Formação Geral da questão 1.
Observa -se que a maior frequência ( exatos 25 \% ) corresponde à dos estudantes que obtiveram nota no intervalo ( 10 ; 20 ], seguidos de 14,7 \% de estudantes alocados igualmente nos intervalos ( 0 ; 10 ] e ( 40 ; 50 ].
Destacam -se, também, os estudantes que deixaram ambas a questão em branco, representando 5,8 \% do total. 6.3.1.7 Comentários sobre a correção de Língua Portuguesa das respostas à Questão Discursiva 1 A Questão 1 se contextualizava em \textbf{torno} do tema"desastres ambientais", tomando como base dois fragmentos de textos acadêmicos publicados em sites oficiais de instituições de ensino superior na internet.
O texto motivador 1 cita um levantamento da ocorrência de desastres naturais associados a fenômenos naturais no Brasil e o texto motivador 2 cita um relatório da ONU sobre a necessidade de minimizar os riscos e os impactos dos futuros desastres naturais.
Com base nessa reflexão inicial, o enunciado solicitava que o participante apresentasse duas propostas de intervenção no âmbito da sustentabilidade socioambiental, de modo a contemplar ações de restauração ou recuperação após a ocorrência desses desastres.
Apesar de não explicitar a exigência de elaboração de um texto dissertativo, o enunciado detalhava várias informações que deveriam compor as duas propostas a serem apresentadas na resposta.
Para atender a essas solicitações, o participante deveria, para \textbf{cada} proposta, elaborar um texto expositivo para detalhar o tipo de desastre ambiental a que \textbf{estava} se referindo, justificar a proposta de intervenção em função do objetivo de garantir a sustentabilidade ambiental e descrever as ações de restauração ou recuperação inseridas na proposta.
Para atender a essas exigências, o vocabulário utilizado deveria ter algum perfil mais técnico e manter coerência com a situação envolvida.
A análise do desempenho linguístico esperado de um estudante que \textbf{está} concluindo seu curso superior apoia -se no princípio de que ele deve utilizar seus conhecimentos sobre o tema para estruturar seus textos de acordo com as características formais da situação comunicativa.
Essa configuração determina exigências quanto à adequação da seleção vocabular, ao desenvolvimento coerente do conteúdo, à estruturação sintática dos períodos, à organização lógica das ideias, à utilização de procedimentos de encadeamento textual e de referenciação, à obediência às exigências morfossintáticas próprias da modalidade escrita da norma-padrão, ao respeito às regras ortográficas e às regras de acentuação gráfica.
Assim, tomando como parâmetro o que pode ser considerado como um desempenho linguístico ideal, a avaliação das respostas procurou mapear detalhadamente o domínio dos participantes quanto aos recursos disponíveis na Língua Portuguesa em situações de escrita formal.
Dessa forma, levaram -se em consideração, no padrão de resposta, as seguintes competências linguísticas próprias da escrita formal culta, responsáveis pelas condições de textualidade : domínio das convenções ortográficas, domínio dos procedimentos de estruturação textual, domínio das convenções de caráter morfossintático e domínio da seleção vocabular.
Essas quatro competências foram reunidas em três itens de avaliação, atendendo aos aspectos explicitados anteriormente : - domínio das convenções ortográficas : grafia de vogais e consoantes, uso de maiúsculas e minúsculas, emprego do hífen e acentuação gráfica ; - domínio dos procedimentos de estruturação textual do ponto de vista microestrutural : organização interna dos períodos, emprego de conectores para a articulação lógica entre os períodos e entre os parágrafos ; emprego de marcas de referenciação lexical e pronominal ; utilização dos sinais de pontuação que contribuem para a organização lógica da frase ; paragrafação do texto ; - domínio das convenções de caráter morfossintático estabelecidas como modelares do ponto de vista da modalidade escrita formal culta da Língua Portuguesa : concordância nominal e verbal, regência nominal e verbal, colocação pronominal, flexão nominal e verbal, correlação entre tempos e modos verbais, ausência de marcas de oralidade.
O domínio da seleção vocabular foi incorporado a essa última competência, tendo em vista a intersecção entre as duas do ponto de vista das exigências do registro formal da modalidade escrita culta.
Os critérios para análise de \textbf{cada} uma das competências que envolvem o domínio das convenções ortográficas, da estruturação textual, das convenções morfossintáticas e da seleção vocabular da escrita formal culta da Língua Portuguesa vêm explicitados a seguir.
O padrão de resposta foi aplicado integralmente para a avaliação das respostas da Questão 1 ( aspectos ortográficos, textuais, morfossintáticos e vocabulares ).
Apesar de o enunciado não fazer referência à exigência de elaboração de um texto dissertativo sobre o tema, a grande maioria dos participantes procurou desenvolver as duas propostas de intervenção de modo detalhado, o que permitiu a avaliação dos quatro aspectos linguísticos do padrão de resposta.
Observou -se, entretanto, uma significativa diversidade de situações de resposta desde a citação de duas propostas de ação sob a forma de frases nominais ( exemplo :"esclarecimento à população sobre os riscos ambientais") ou verbais ( exemplo : Aspectos morfossintáticos Esta competência envolve :  concordância nominal e verbal ;  regência nominal e verbal ;  flexão nominal e verbal ;  correlação entre modos e tempos verbais no período ;  colocação pronominal.
Espera -se que o participante :  flexione o verbo para estabelecer concordância de número com o sujeito da frase ;  flexione o artigo, o adjetivo e o pronome para concordar em número e em gênero com o substantivo a que se referem ;  observe a regência nominal e a verbal, utilizando a preposição adequada depois de um substantivo, um verbo ou um adjetivo ;  empregue adequadamente o acento grave indicador da crase entre uma preposição e um artigo ( a+a ) ;  obedeça às regras de colocação pronominal ( próclise e ênclise ), distintas dos hábitos da oralidade ou da escrita informal ;  flexione adequadamente verbos, substantivos, adjetivos e pronomes no que diz respeito à expressão das categorias gramaticais ;  flexione os verbos para expressar a correlação de modo e tempo nas estruturas subordinadas.
Aspectos vocabulares Esta competência envolve :  precisão na utilização do vocabulário relacionado ao tema da questão ;  ausência de marcas de oralidade, como termos de sentido muito genérico e termos de registros mais informais.
Espera -se que o participante :  utilize repertório vocabular diversificado, sem repetição inadequada de termos ;  recorra a vocabulário preciso em relação ao tema da questão ;  complemente o sentido das palavras por meio de expressões modificadoras ( exemplo : “ Desenvolvimento de projetos ” – sentido incompleto no que diz respeito à palavra “ desenvolvimento ” e à palavra “ projetos ” ) ;  empregue um vocabulário mais formal, com ausência de marcas da oralidade, como termos de sentido muito genérico ( “ coisa ”, “ negócio ”, “ você ” ) e termos de registros mais informais ( como gírias, jargões, frases feitas, ditados populares, termos regionais ) ;  seja coerente na combinação vocabular.
O padrão de resposta foi aplicado integralmente para a avaliação das respostas da Questão 1 ( aspectos ortográficos, textuais, morfossintáticos e vocabulares ).
Apesar de o enunciado não fazer referência à exigência de elaboração de um texto dissertativo sobre o tema, a grande maioria dos participantes procurou desenvolver as duas propostas de intervenção de modo detalhado, o que permitiu a avaliação dos quatro aspectos linguísticos do padrão de resposta.
Observou -se, entretanto, uma significativa diversidade de situações de resposta desde a citação de duas propostas de ação sob a forma de frases nominais ( exemplo :"esclarecimento à população sobre os riscos ambientais") ou verbais ( exemplo :"evitar a propagação do efeito estufa"), sem desenvolvimento ou justificativas, até a construção de respostas com \textbf{detalhamento} das propostas apresentadas, acompanhadas por"evitar a propagação do efeito estufa"), sem desenvolvimento ou justificativas, até a construção de respostas com \textbf{detalhamento} das propostas apresentadas, acompanhadas por justificativas.
Com base na avaliação das respostas da Questão 1, pode -se considerar que o desempenho dos participantes foi de nível mediano, em sua maior parte, com desvios pontuais nas quatro competências analisadas.
As respostas que foram avaliadas como fracas são de dois tipos principais : respostas itemizadas, com uma frase curta para \textbf{cada} proposta apresentada ; ou respostas constituídas por pequenos textos, com muitos desvios nas quatro competências, principalmente desestruturação sintática e falta de domínio das convenções morfossintáticas.
Nas respostas que foram avaliadas como medianas, os estudantes apresentaram uma estruturação textual relativamente adequada, com ocorrências pontuais de desvios nas quatro competências.
Entretanto, em todos os aspectos, também houve estruturas em que se respeitaram as características da modalidade escrita formal do padrão culto da Língua Portuguesa.
Nas respostas que foram avaliadas como boas, os estudantes apresentaram desenvolvimento textual fluente e coerente, vocabulário formal e técnico, domínio das convenções morfossintáticas e ortográficas.
Na sequência, para \textbf{cada} uma das competências analisadas durante a correção das respostas dos participantes para a Questão 1, destacam -se as principais características observadas.
Aspectos ortográficos – O desempenho dos participantes revelou uma diferença muito grande nos dois aspectos analisados, como tem sido observado em edições anteriores do Enade : baixo índice de desvios da grafia padrão e grande índice de desvios de acentuação.
Em vários casos, ocorre ausência quase completa de acentuação gráfica.
Com base na análise das respostas, conclui -se, portanto, que, de uma maneira geral, os participantes dominam as convenções relativas à grafia das palavras, mas desconhecem, ou desprezam, as convenções relativas à acentuação gráfica.
Seguem exemplos de algumas ocorrências : a ) eliminação do acento indicador da sílaba tônica : - palavras proparoxítonas : “ tecnica ” ( por"técnica"), “ proposito ” ( por propósito"), “ catastrofe ” ( por catástrofe"), “ ambito ” ( por"âmbito"),"sustentavel"( por"sustentável") ; b ) palavras paroxítonas terminadas em ditongo crescente : “ reincidencia ” ( por"reincidência"),"industrias"( por"indústrias") ; c ) palavras oxítonas : “ comite ” ( por"comitê") ; b ) uso indevido do acento gráfico, localizado em sílaba átona : “ ecônomia ” ( por"economia"), “ prevênção ” ( por"prevenção"), “ intervênção ” ( por intervenção"),"reciclavél"( por"reciclável") ; c ) omissão da cedilha, do til e do traço do T :"seguranca"( por"segurança"),"realizacao"( por"realização"),"proibicao"( por"proibição").
Há, inclusive, respostas com ausência total de cedilha ou til ; d ) desvios de grafia :"concientizar"( por"conscientizar"), “ atravéz ” ( por"através"), “ intrui ” ( por “ instrui ” ), “ intruídos ” ( por “ instruídos ” ), “ extrangeiro ” ( por"estrangeiro"), “ envestir ” ( por"investir"), “ tecnoligia ” ( por"tecnologia"), “ evulução ” ( por"evolução"), “ consiliar ” ( por"conciliar"), “ extrutura ” ( por"estrutura") ;"serto ” ( por"\textbf{certo}"),"siguinificativos ” ( por"significativos"), “ esportamos ” ( por"exportamos") ; e ) omissão de sílaba : “ natuza ” por “ natureza ” ; f ) desvios de segmentação : “ apartir ” ( por"a partir"), “ a cerca ” ( por “ acerca ” ), “ oque ” ( por"o que") ; g ) uso indevido de inicial maiúscula :"Boletos"; h ) ao contrário do que se esperava, apareceram raramente as abreviaturas próprias do “ internetês ”, relacionadas ao uso de redes sociais e emails, como"p/ ”, “ \textbf{tá} ”, “ pra ”, “ pro ”, “ prum ”.
Aspectos morfossintáticos – De uma maneira geral, pode -se afirmar que há um domínio mediano das convenções morfossintáticas pelos participantes.
A seguir analisam -se alguns desvios relativos à regência e à concordância nominal e verbal.
Quanto à regência nominal e verbal, o desvio mais frequente é a falta do sinal indicativo da crase, o que revela que o usuário não tem consciência de que “ a ” reúne um artigo e uma preposição “ a ”, exigida pela regência do termo anterior ( substantivo, adjetivo ou verbo ).
Outro problema relacionado à regência verbal e à nominal é a ausência de preposição antes de pronome relativo, processo generalizado na modalidade oral da língua, em situações de registro informal.
Também é recorrente o emprego inadequado de preposição, por exemplo : “ ações a evitá -los ”.
Quanto à concordância nominal e verbal, seis aspectos merecem destaque : - palavras paroxítonas terminadas em ditongo crescente : “ reincidencia ” ( por"reincidência"),"industrias"( por"indústrias") ; - palavras oxítonas : “ comite ” ( por"comitê") ; b ) uso indevido do acento gráfico, localizado em sílaba átona : “ ecônomia ” ( por"economia"), “ prevênção ” ( por"prevenção"), “ intervênção ” ( por intervenção"),"reciclavél"( por"reciclável") ; c ) omissão da cedilha, do til e do traço do T :"seguranca"( por"segurança"),"realizacao"( por"realização"),"proibicao"( por"proibição").
Há, inclusive, respostas com ausência total de cedilha ou til ; d ) desvios de grafia :"concientizar"( por"conscientizar"), “ atravéz ” ( por"através"), “ intrui ” ( por “ instrui ” ), “ intruídos ” ( por “ instruídos ” ), “ extrangeiro ” ( por"estrangeiro"), “ envestir ” ( por"investir"), “ tecnoligia ” ( por"tecnologia"), “ evulução ” ( por"evolução"), “ consiliar ” ( por"conciliar"), “ extrutura ” ( por"estrutura") ;"serto ” ( por"\textbf{certo}"),"siguinificativos ” ( por"significativos"), “ esportamos ” ( por"exportamos") ; e ) omissão de sílaba : “ natuza ” por “ natureza ” ; f ) desvios de segmentação : “ apartir ” ( por"a partir"), “ a cerca ” ( por “ acerca ” ), “ oque ” ( por"o que") ; g ) uso indevido de inicial maiúscula :"Boletos"; h ) ao contrário do que se esperava, apareceram raramente as abreviaturas próprias do “ internetês ”, relacionadas ao uso de redes sociais e emails, como"p/ ”, “ \textbf{tá} ”, “ pra ”, “ pro ”, “ prum ”.
Aspectos morfossintáticos – De uma maneira geral, pode -se afirmar que há um domínio mediano das convenções morfossintáticas pelos participantes.
A seguir analisam -se alguns desvios relativos à regência e à concordância nominal e verbal.
Quanto à regência nominal e verbal, o desvio mais frequente é a falta do sinal indicativo da crase, o que revela que o usuário não tem consciência de que “ a ” reúne um artigo e uma preposição “ a ”, exigida pela regência do termo anterior ( substantivo, adjetivo ou verbo ).
Outro problema relacionado à regência verbal e à nominal é a ausência de preposição antes de pronome relativo, processo generalizado na modalidade oral da língua, em situações de registro informal.
Também é recorrente o emprego inadequado de preposição, por exemplo : “ ações a evitá -los ”.
Quanto à concordância nominal e verbal, seis aspectos merecem destaque : a ) ausência ou presença indevida de marca de plural ( com sujeito anteposto ou posposto ).
Exemplos : “ para que a distribuição de água cheguem ”, “ fica claro os benefícios ” ;"evitar que aconteça esses desastres";"criar um sistema de coletagem de ideias para que seja amenizada as perdas";"equipamentos que medi ” ; “ Dessa forma corroborando para que esferas do nível federal enalteça mediadas a serem tomadas... ” ;"para que houvessem mais áreas drenadas"; b ) aparecimento da marca de plural em verbos ou adjetivos relacionados a núcleos substantivos no singular, devido à presença de modificadores no plural.
Exemplos :"o aumento dos acidentes ambientais devem ser evitados ” ; “ a fiscalização sobre os laudos deveriam ”, evidenciando um processo de hipercorreção ; c ) ausência de acento circunflexo na forma plural do presente do indicativo do verbo “ ter ”, que é considerada como um desvio na concordância verbal e não na acentuação gráfica :"os estados e municípios tem que aplicar medidas punitivas a empresas que desmatam"; d ) ausência de concordância de gênero, no âmbito de sintagmas nominais longos, em que o adjetivo antecede ou \textbf{está} afastado do substantivo :"seria necessário uma reeducação";"é necessário a mobilização"; e ) emprego indevido de pronomes, reproduzindo hábito da oralidade, por exemplo, uso do pronome reto como objeto : “ deixar elas isoladas ” ; f ) confusão entre formas verbais : “ apoia -se ” por “ apoiasse".
Várias marcas de oralidade foram identificadas, embora não em alta frequência : o uso do pronome relativo “ onde ” como relativo universal, falta de artigo definido antes de substantivo, expressões informais, eliminação de preposições.
Aspectos Vocabulares – O desempenho nesta competência foi muito diversificado, revelando diferentes graus de domínio vocabular.
Como o Enade/2019 avaliou graduandos de Engenharia, pode -se observar nitidamente, em uma parte dos participantes, o domínio de um vocabulário específico na discussão da questão dos desastres ambientais e das propostas de intervenção para esse grave problema.
Por outro lado, uma grande quantidade de respostas limitou -se ao senso comum, marcado por um vocabulário pouco específico, de caráter genérico.
Quanto ao grau de formalidade do vocabulário empregado, pode -se afirmar que os participantes empregaram um registro semiformal, sem a utilização de gírias ou palavras típicas da oralidade.
Vários tipos de inadequação foram observados na seleção e na utilização do vocabulário de uma parte dos participantes : a ) excessiva repetição de \textbf{certas} palavras, revelando limitação de repertório vocabular ; b ) expressões da oralidade – por exemplo : “ colocar isso na cabeça das pessoas ”, “ não saem do lugar ” ; c ) seleção vocabular incompatível com o contexto, gerando falta de inteligibilidade, por exemplo : “ no requisito ao combate ” ; “ moradias dignas de baixo custo ” ; d ) falta de domínio de vocabulário mais abstrato e de maior complexidade, essencial ao desenvolvimento do texto de base dissertativa : uso de"corroborando"no lugar de"colaborando"; “ suposto ” por “ pressuposto"; “ devemos obter algumas atitudes"; “ recorrentes ” por “ decorrentes"; “ destacados ” por “ descartados"; “ eminente ” por “ iminente ” ; e ) emprego de palavras inexistentes, por exemplo : “ trabalhabilidade ”, “ dizrespeito ” ;"acessiva ”,"coletagem"; f ) imprecisão vocabular, como o uso da palavra “ coisa ” em vários textos, demonstrando o domínio precário da norma-padrão da Língua Portuguesa, com presença constante de traços de oralidade.
Aspectos textuais – A maior parte das respostas foi constituída por textos estruturados, com justificativas para as propostas apresentadas, constituindo um pequeno texto de base dissertativa, mas dividido em dois parágrafos independentes, sem conexão lógica, devido ao que foi solicitado pelo comando.
Entretanto, foram identificadas algumas estruturas textuais que não atenderam à necessidade de se construir um texto para apresentação das propostas de intervenção : a ) respostas com sintagmas nominais curtos, em geral, de duas ou três palavras.
Exemplos :"secas, inundações, alagamentos";"redução dos níveis de risco"; b ) respostas com desenvolvimento do conteúdo da palavra ou expressão nuclear ( sintagmas nominais mais amplos, em geral de mais de quatro palavras ).
Exemplos :"minimização dos impactos futuros dos desastres ambientais";"fortalecimento das resiliências ambiental, social e econômica";"contenção de encostas em áreas de risco"; c ) respostas com estrutura frasal completa, mas com pouco desenvolvimento, reduzidas a apenas um período : Exemplos : “ dragar ” os rios que cortam as grandes cidades";"estimular o trabalho voluntário".
As respostas que se constituíram como textos desenvolvidos apresentaram distintos desempenhos : textos bem elaborados, com estruturação, fluência e coerência ; textos de desempenho mediano, com alguns desvios de estruturação ; e textos bem precários, com muitos desvios e comprometimento da coerência e da fluência textual, refletindo tendências gerais na escrita de estudantes de diferentes níveis de ensino.
Seguem alguns desvios observados : a ) sequência justaposta de ideias sem encaixamentos sintáticos ; b ) redução de estruturas subordinadas, ao lado do aumento na frequência de estruturas coordenadas e absolutas ; c ) redução no uso de conectores para expressar relações lógicas essenciais à construção do texto ; d ) repetição exaustiva de termos sem a utilização de procedimentos mais sofisticados de referenciação ( hiperonímias, hiponímias, nominalizações, expressões metafóricas ) ; e ) frases fragmentadas que comprometem a estrutura lógico-gramatical, como frases formadas apenas por oração subordinada, sem oração principal.
Exemplos : - oração subordinada solta : “ No qual forneceria moradias dignas de baixo custo ” ; - frase iniciada por gerúndio solta : “ Evitando assim a ocupação de lugares impropios ( por impróprios ) para viver ” ;"Podendo também fazer eventos na intenção..."; - truncamento sintático : “... isso faz com que os bueiros fiquem entupidos e aconteça e a água...".
Quanto à utilização dos sinais de pontuação, observou -se uma grande diversidade de situações nos textos analisados, desde textos bem pontuados até textos completamente sem pontuação, inclusive sem ponto final.
São os seguintes os tipos de problemas encontrados : - emprego de vírgula para separar o sujeito e o predicado, procedimento muito frequente ; - emprego de vírgula no lugar do ponto para separar ideias que constituem períodos distintos ; - ocorrência de apenas uma das vírgulas para destacar uma palavra, uma expressão ou uma oração encaixada ; - ausência de vírgula para separar elementos de uma enumeração ; - ausência de vírgula para separar oração adjetiva explicativa ou utilização inadequada de vírgula para separar oração adjetiva restritiva ; - emprego de ponto e vírgula no lugar de vírgula ; - ausência de ponto final para encerrar o texto. 6.3.1.8 Análise de Língua Portuguesa da Questão Discursiva 2 do Componente de Formação Geral Os dados de Engenharia Florestal, obtidos a partir das respostas à questão discursiva 2 do Componente de Formação Geral, no que tange à Língua Portuguesa, encontram -se na Tabela 6.17 e no Gráfico 6.10.
Nesse aspecto, os estudantes de todo o Brasil obtiveram Média 54,6.
A maior Média em relação à Língua Portuguesa foi obtida na região Sudeste ( 60,2 ), e a menor, na região Norte ( 50,4 ).
Quanto à variabilidade das notas, o Desvio padrão de todo o Brasil foi 22,2.
O menor Desvio padrão foi obtido na região Sudeste ( 18,7 ), e o maior Desvio padrão foi obtido na região Norte ( 24,7 ).
A Mediana das notas de Língua Portuguesa para a questão 2 foi 55,0 para o Brasil, como um todo, a mesma obtida nas regiões Norte, Nordeste e Centro-Oeste.
A Mediana das regiões Sudeste e Sul foi 60,0.
A nota Máxima para todo o Brasil foi 90,0, com, pelo menos, um estudante tirando essa nota em todas as regiões, sem exceção.
Já a nota Mínima foi zero em todas as regiões do país.
No Gráfico 6.10, mostra -se a distribuição das notas de Língua Portuguesa do Componente de Formação Geral.
Observa -se que a maior frequência ( 26,4 \% ) corresponde à dos estudantes que obtiveram nota no intervalo ( 50 ; 60 ], seguido por 24,5 \% dos estudantes alocados no intervalo ( 60 ; 70 ].
Destacam -se, também, os estudantes que deixaram a questão em branco, representando 7,6 \% do total. 6.3.1.9 Comentários sobre a correção de Língua Portuguesa das respostas à Questão Discursiva 2 A Questão 2 de Formação Geral abordava o tema do desenvolvimento científico e do desenvolvimento tecnológico do Brasil, tomando como base um fragmento de texto publicado em site oficial de um órgão governamental de pesquisa na internet, o Ipea.
O texto motivador informava que o Brasil \textbf{está} em posição intermediária em todos os indicadores de produção e utilização de conhecimentos e novas tecnologias.
Além desse texto, a questão apresentava uma imagem que estabelecia uma ligação entre Instituição de Ensino Superior - Empresa - Setor Público.
Com base nessa reflexão inicial despertada pelo texto motivador, o comando solicitava que o participante citasse dois ganhos para o campo científico e dois ganhos para o campo econômico no país, resultantes da relação entre Instituição de Ensino Superior - Empresa - Setor Público.
Como o comando solicitava apenas uma “ citação ” de ganhos, sem referência à exigência de justificativas ou explicações que ensejariam a elaboração de um texto dissertativo sobre o tema, observou -se uma grande fragmentação do pensamento abstrato e, por consequência, uma grande fragmentação discursiva.
Essa diversidade pode ter sido consequência de uma \textbf{certa} insegurança entre os participantes a respeito de qual modelo de resposta produzir.
Assim, houve uma grande diversidade de modelos de resposta, desde a citação de duas palavras ou expressões, sem desenvolvimento, até a construção de respostas com \textbf{detalhamento} dos ganhos sugeridos, acompanhados por uma justificativa de caráter dissertativo, passando por vários modelos intermediários.
Essa diversidade tipológica, que será mais detalhada adiante, exigiu mudanças nos critérios habitualmente utilizados para avaliação de uma produção textual.
Como resultado desse enunciado, as respostas foram majoritariamente apresentadas sob a forma de dois itens, \textbf{cada} um atendendo à"citação"dos ganhos científicos e dos ganhos econômicos.
O que variou foi o formato textual dessa"citação", conforme detalhado a seguir. a ) Respostas com palavras ou expressões ( sintagmas nominais curtos, em geral, de uma a três palavras ), constituindo uma frase nominal com núcleo substantivo : Exemplos : “ Descobertas de algumas doenças ” ; “ Frigoríficos de frango ” ; “ Valorização ” ; “ Infraestrutura ” ; “ Laboratórios ” ; “ Campos de estudo ” ; “ Tecnologia ” ; “ Automação ”. b ) Respostas com desenvolvimento do conteúdo da palavra ou expressão nuclear ( sintagmas nominais mais amplos, em geral de mais de quatro palavras ), constituindo uma frase nominal com núcleo substantivo : Exemplos : “ Aquisição de equipamentos e materiais necessários ao desenvolvimento de pesquisas ” ; “ Descobertas no ramo da medicina relacionadas à cura de novas doenças ” ; “ Descoberta de novas tecnologias pelas instituições de pesquisa ” ; “ Melhoria das instituições de ensino superior pela ajuda empresarial ” ; “ Desenvolvimento de energias sustentáveis e economicamente acessíveis ”. c ) Respostas iniciadas por palavra ou expressão ( sintagmas nominais curtos ), mas com a presença de expressão verbal na complementação do núcleo sintático, constituindo uma frase nominal ( com núcleo substantivo ) complexa devido à presença de oração subordinada com núcleo verbal : Exemplos : “ Setor econômico possibilitando a oferta de emprego ao cidadão ” ; “ Auxílio financeiro para que os pesquisadores possam permanecer na academia ” ; “ Desenvolvimento de novas metodologias que permitam aumentar a produtividade de áreas agriculturáveis brasileiras ” ; “ Flexibilização de leis ( ou criação ) de modo a estimular a utilização de novas tecnologias ” ; “ Maior investimento estrangeiro para incentivar novas pesquisas e reconhecimento internacional ”. d ) Respostas com estrutura frasal completa ( com núcleo verbal ), mas com pouco desenvolvimento, reduzidas a apenas um período : Exemplos : “ Possibilitar a oferta de emprego ao cidadão e investimentos para as empresas, com fácil acesso para o pequeno empresário ” ; “ Investindo nesses recursos de pesquisas cientificas, obteve -se um retorno na economia do país ” e ) Respostas mais desenvolvidas, com justificativas, constituindo um pequeno texto de base dissertativa, dividido ou não em dois parágrafos independentes, sem conexão lógica : Exemplo : “ A consultoria de membros laboratoriais de instituições de ensino superior em pequenas e médias empresas pode agregar valor tanto no processo produtivo da empresa quanto numa maior qualificação desses “ consultores ” no mercado de trabalho. ( Maior empregabilidade ) ”.
Para avaliar essa diversidade de situações discursivas da Questão 2, não foi possível utilizar o padrão de resposta que vem sendo utilizado para a avaliação do desempenho linguístico dos participantes dos Enade anteriores e na correção da Questão 1.
Aquele padrão se revelou inadequado para a realidade das respostas analisadas na Questão 2, porque uma parcela significativa dos participantes elaborou respostas curtas, compostas por itens isolados, criando praticamente uma lista de palavras, expressões ou frases soltas, sem conexão semântica ou textual.
Assim, realizou -se uma adequação dos critérios de avaliação para atender às características das respostas produzidas, em função das exigências estabelecidas pelo enunciado : “ citar ganhos ”.
A análise do desempenho linguístico esperado de um estudante que \textbf{está} concluindo seu curso superior e, devido a isso, \textbf{está} realizando o exame do Enade, tem se pautado no princípio de que ele deve utilizar seus conhecimentos sobre os dois temas de Formação Geral apresentados e estruturar seus textos, de base dissertativa, de acordo com as características do registro formal adequado à situação comunicativa – avaliação de conhecimentos.
Essa configuração determina exigências quanto : à adequação da seleção vocabular, à progressão temática coerente do conteúdo, à estruturação sintática dos períodos, à organização lógica das ideias, à utilização de procedimentos de encadeamento textual e de referenciação, à obediência às exigências morfossintáticas próprias da modalidade escrita da norma-padrão, ao respeito às regras ortográficas e às regras de acentuação gráfica.
São competências distintas, que permitem um mapeamento detalhado do domínio dos recursos disponíveis na Língua Portuguesa para a comunicação escrita formal, que constituem um acervo linguístico dos graduandos para enfrentar as exigências de sua vida profissional futura.
Na Questão 2 da edição do Enade de 2019, entretanto, a avaliação dos aspectos textuais ficou extremamente prejudicada, porque a maioria das respostas não apresentou material linguístico suficiente para que essa competência fosse pontuada.
Devido a isso, a competência relativa ao domínio dos aspectos textuais não foi avaliada na Questão 2 do Enade/2019 e, em compensação, foram ampliados, consequentemente, os valores relativos às demais competências : atendimento às convenções ortográficas, adequação na seleção vocabular e atendimento às convenções morfossintáticas.
Essa solução se justifica, também, porque a avaliação dos aspectos textuais criaria uma desigualdade ente os participantes : aqueles que atenderam ao comando do enunciado e produziram uma resposta composta apenas por itens para apresentar os “ ganhos ” seriam prejudicados, porque seriam penalizados por não atenderem às exigências de uma construção coerente e bem desenvolvida.
O desempenho linguístico dos participantes foi avaliado, portanto, com base nas três competências descritas a seguir : - domínio das convenções ortográficas : grafia de vogais e consoantes, uso de maiúsculas e minúsculas, emprego do hífen e acentuação gráfica ; - domínio da seleção vocabular, tendo em vista as exigências semânticas do tema desenvolvido e do registro formal da modalidade escrita culta ; - domínio das convenções de caráter morfossintático estabelecidas como modelares do ponto de vista da modalidade escrita formal culta da Língua Portuguesa : concordância nominal e verbal, regência nominal e verbal, colocação pronominal, flexão nominal e verbal, correlação entre tempos e modos verbais, ausência de marcas de oralidade.
Os critérios para análise de \textbf{cada} uma das competências que envolvem o domínio das convenções ortográficas, das convenções morfossintáticas e da seleção vocabular da escrita formal culta da Língua Portuguesa vêm explicitados a seguir.
O padrão de resposta aplicado para a avaliação da Questão 2 não considerou, portanto, as competências relativas aos aspectos textuais, conforme foi realizado para a Questão 1.
Ao final do processo de avaliação, ficou bem evidenciada a influência que o enunciado da Questão 2 exerceu sobre o desempenho dos participantes.
A grande diversidade de situações de resposta não expressou apenas a dificuldade de organização do pensamento sob a forma de um texto bem estruturado e coerente, mas se refletiu também nas competências relativas aos aspectos vocabulares e morfossintáticos.
Por exemplo, nas respostas mais sintéticas, constituídas por itens preenchidos por sintagmas nominais simples ( compostos apenas por um núcleo substantivo ), os estudantes apresentaram fraco desempenho nos aspectos vocabulares ao empregarem termos muito vagos ou imprecisos semanticamente ( exemplo : “ Hierarquia ” ).
Quando o núcleo desse modelo de resposta era preenchido por substantivos deverbais ( relacionados semanticamente a um verbo ), ficavam incompletos do ponto de vista morfossintático pela ausência de um complemento nominal ( exemplo : “ Valorização ” de quê? ).
Com base na avaliação das respostas, pode -se considerar que o desempenho dos participantes foi de nível mediano, em sua maior parte, com vários desvios nas três competências analisadas.
As respostas que foram avaliadas como fracas foram de dois tipos principais : respostas itemizadas, com sintagmas reduzidos a um núcleo, acompanhado ou não por modificadores ; ou respostas itemizadas constituídas por pequenos parágrafos, com muitos desvios nas três competências, principalmente na seleção vocabular e no domínio das convenções morfossintáticas.
Em comum nos dois tipos, a pouca ou nenhuma articulação entre os itens ‘ a ’ e ‘ b ’, propostos na questão.
As respostas que foram avaliadas como medianas também foram itemizadas, mas \textbf{cada} item era abordado por parágrafos com uma estruturação textual relativamente adequada e desvios nas três competências.
As respostas que foram avaliadas como boas apresentavam desenvolvimento textual fluente e coerente, vocabulário formal e técnico, domínio das convenções morfossintáticas e ortográficas, mesmo que apresentassem eventuais desvios em um ou mais dos aspectos avaliados.
Apresentam -se, a seguir, algumas observações sobre o desempenho dos participantes em relação às três competências avaliadas.
Aspectos ortográficos – Assim como foi observado na Questão 1, o desempenho dos participantes revelou uma diferença muito grande nos dois aspectos analisados : baixo índice de desvios de grafia padrão e grande índice de desvios de acentuação.
Em vários casos, ocorre ausência quase completa de acentuação gráfica.
Com base na análise das respostas, conclui -se, portanto, que, de uma maneira geral, os participantes dominam as convenções relativas à grafia das palavras, mas desconhecem, ou desprezam, as convenções relativas à acentuação gráfica.
Seguem exemplos de algumas ocorrências : a ) eliminação do acento indicador da sílaba tônica : - palavras proparoxítonas : “ praticos ” ( por “ práticos ” ), “ tecnologicos ” ( por “ tecnológicos ” ) ; - palavras paroxítonas : “ diminuissem ” ( por “ diminuíssem ” ), “ relevancia ” ( por “ relevância ” ) ; “ inimaginaveis ” ( por “ inimagináveis ” ) ; b ) uso indevido do acento gráfico, localizado em sílaba átona : “ cancêr ” ( por “ câncer ) ; c ) omissão da cedilha, do til ou do traço da letra"t": “ populacão ” ( por “ população ” ), “ especulaçoes ” ( por “ especulações ” ), “ orientacoes ” ( por “ orientações ” ) ; d ) desvios de grafia :"concientizar"( por"conscientizar"), “ tramisiveis ” ( por “ transmissíveis ” ), “ pulblico ” ( por “ público ” ) ; e ) uso indevido de inicial maiúscula : “ Periódicos ” ; f ) uso indevido ou ausência de hífen : “ micro-empresário ” ( por “ microempresário ” ), “ recém formados ” ( por “ recém-formados ” ) ; g ) ao contrário do que se esperava, foram poucas as ocorrências de abreviaturas próprias do “ internetês ”, relacionadas ao uso de redes sociais e emails, como"p/ ”, “ \textbf{tá} ”, “ pra ”, “ pro ”, “ prum ”.
Aspectos Vocabulares – O desempenho nesta competência foi muito diversificado, revelando diferentes graus de domínio vocabular.
A seleção e a adequação vocabular foram consideradas essenciais na avaliação das respostas itemizadas, porque a tendência dos participantes foi a de utilizar palavras muito vagas e genéricas.
Quanto ao grau de formalidade do vocabulário empregado, pode -se afirmar que os participantes utilizaram um registro semiformal, sem a utilização de gírias ou palavras típicas da oralidade.
Vários tipos de inadequação foram observados na seleção e na utilização do vocabulário de uma parte dos participantes : a ) desvios na precisão vocabular : “ Começar a produzir ”, sem a delimitação do que ‘ produzir ” ; “ o custo seria inferior ”, sem a delimitação do que tem custo ; “ Com isso o custo seria inferior ao comprador, onde iria igualar aos países mais Desenvolvidos ”, sem delimitar o que se igualaria aos países mais desenvolvidos ; b ) seleção vocabular inadequada : “ Estudantes que \textbf{estão} \textbf{cada} vez mais inteligentes e com grandes ideias afim de ajudar essa empresa ” ; “ recursos de supra importância ” ;"agregando valor ao nome brasileiro ao exterior ” ; “ os profissionais qualificados através de instituições superiores, promovem desempenho econômico abrangente à economia do país, pois a geração de empregos e presença de concursos públicos privados, geram saldo positivo em contratações e giro da economia promovida pelas empresa do setor público ” ; c ) neologismos não autorizados pelos dicionários : “ tecnologias sustentais ” ; “ Ingressão no meio empresarial ”.
Aspectos morfossintáticos – De uma maneira geral, pode -se afirmar que há um domínio mediano das convenções morfossintáticas pelos participantes, embora a estrutura itemizada das respostas, ocorrida na Questão 2, tenha favorecido muitos desvios morfossintáticos : a ) falta de paralelismo entre elementos que introduzem respostas itemizadas : “ maior número de doutores formados pelas universidades ” ; “ maior reconhecimento mundial relacionado a educação ” ; “ maior investimento de multinacionais, gerando empregos ” ; “ estimular o empreendimento interno ” ( no lugar de “ estímulo ” para manter o paralelismo com os outros substantivos dos itens anteriores ). b ) falta de preposição para manter o paralelismo em enumeração : “ maior geração de conhecimento e ( de ) informação ” ; “ maior disponibilização de recursos e ( de ) dados ” ; “ relatórios de realização e ( de ) atividade econômica ” ; “ a produção qualificada da ciência e ( da ) tecnologia ” ; “ ajudar na qualidade e ( na ) rapidez da produção ”. c ) falta de artigo na enumeração de itens : “ a educação e ( o ) conhecimento elevam uma sociedade ” ; “ onde é melhor realizar melhorias para as cidades, ( os ) \textbf{estados} e ( o ) país ” ; “ isso volta como serviço para as empresas e ( o ) setor público ” ; “ juntamente com o setor privado e ( o ) público ” ; “ criação de novas tecnologias para as empresas e ( o ) setor público ” ; “ desenvolver parceria entre as universidades e ( as ) empresas privadas ” ; “ investimentos viabilizam ( a ) produção e ( a ) venda de tecnologia ” ; “ profissionais capacitados para os trabalhos e ( as ) pesquisas ” ; “ com a contribuição e ( a ) colaboração dos setores públicos e empresariais ” ; “ surgimento das novas ideias, ( das ) pesquisas e ( dos ) projetos ” ; “ o preço cobrado pelos produtos e ( pelas ) empresas seria menor ” ; “ novas descobertas nos campos da saúde e ( da ) engenharia ”. d ) ausência de complementos nominais junto a substantivos abstratos deverbais, principalmente nas respostas itemizadas : “ valorização ”, no lugar de “ valorização de... ” ; “ patrocínio ” no lugar de “ patrocínio de... ”.
Desvios de regência nominal e verbal :  o desvio mais frequente é a falta do sinal indicativo da crase, o que revela que o usuário não tem consciência de que o termo “ a ” reúne um artigo"a"e uma preposição “ a ”, exigida pela regência do termo anterior ( substantivo, adjetivo ou verbo ) : “ maior reconhecimento mundial relacionado a ( por à ) educação ” ;  ausência de preposição antes de pronome relativo, processo generalizado na modalidade oral da língua, em situações de registro informal : “ as instituições de ensino superior ("em") que as empresas investem ” ;  emprego inadequado de preposição : “ o custo seria inferior ao ( por"para o") comprador ”.
Desvios de concordância nominal e verbal :  ausência ou presença indevida de marca de plural ( com sujeito anteposto ou posposto ).
Exemplos : “ Os setores público e ensino superior poderiam desenvolverem ( por “ desenvolver ” ) pesquisas para empresas seja ( por “ sejam ” ) elas do tipo tecnologico ( por “ tecnológico ” ), saúde, transporte entre outros ” ; “ já as empresas poderiam financiar projetos de pesquisas para os setores público e ensino superior que também poderia se utilizarem ( por “ utilizar ” ) de suas partes estruturais e acadêmicas para desenvolvimento tecnologico ” ( por “ tecnológico ” ) ; “ as empresas que produz ( por “ produzem ” ) os materiais ” ; “ o engajamento entre ambos elevam ( por “ eleva ” ) a economia ”. 6.3.1.10 Considerações finais sobre a correção de Língua Portuguesa das respostas às Questões Discursivas do Componente de Formação Geral A análise do desempenho linguístico das duas questões de Formação Geral do Enade/2019 permite chegar a algumas conclusões que podem colaborar para um aperfeiçoamento dos cursos superiores e para uma adequação do exame aos seus propósitos institucionais.
Um desses propósitos é a avaliação sobre a competência discursiva dos graduandos, conforme determinação do Inep : as questões discursivas devem fornecer material para que as respostas sejam avaliadas em aspectos como clareza, coerência, coesão, estratégias argumentativas, utilização de vocabulário adequado e correção gramatical do texto.
Com base na constatação de que nossos estudantes têm um domínio apenas mediano dessas exigências de construção de um texto formal culto, conclui -se que é preciso  Aparecimento da marca de plural em verbos ou adjetivos relacionados a núcleos substantivos no singular devido à presença de modificadores no plural.
Exemplos :"o financiamento das bolsas de estudo devem ( por “ deve ” ) ser ampliados ” ( por"ampliado ” ) ; “ o acesso aos cursos deveriam ( por “ deveria ” ) ser incentivados ( por"incentivado") pelas empresas ”, evidenciando um processo de hipercorreção.  Ausência de acento circunflexo na forma plural do presente do indicativo dos verbos “ ter ” e “ vir ”, considerada como um desvio na concordância verbal e não na acentuação gráfica.
Exemplos : “ ganhos na economia vem ( por “ vêm ” ) através dos aumentos com exportação ” ;"as universidades públicas tem ( por"têm") recebido menor quantidade de verbas para suas atividades".  Ausência de concordância de gênero, no âmbito de sintagmas nominais longos, em que o adjetivo antecede ou \textbf{está} afastado do substantivo.
Exemplos :"seria necessário uma reeducação"( por"necessária") ;"é necessário ( por"necessária") a mobilização";  Emprego de “ onde ” não locativo.
Exemplo : “ Com isso o custo seria inferior ao comprador, onde iria igualar aos países mais desenvolvidos ”. 6.3.1.10 Considerações finais sobre a correção de Língua Portuguesa das respostas às Questões Discursivas do Componente de Formação Geral A análise do desempenho linguístico das duas questões de Formação Geral do Enade/2019 permite chegar a algumas conclusões que podem colaborar para um aperfeiçoamento dos cursos superiores e para uma adequação do exame aos seus propósitos institucionais.
Um desses propósitos é a avaliação sobre a competência discursiva dos graduandos, conforme determinação do Inep : as questões discursivas devem fornecer material para que as respostas sejam avaliadas em aspectos como clareza, coerência, coesão, estratégias argumentativas, utilização de vocabulário adequado e correção gramatical do texto.
Com base na constatação de que nossos estudantes têm um domínio apenas mediano dessas exigências de construção de um texto formal culto, conclui -se que é preciso aperfeiçoar sua formação superior, dando -lhes acesso às ferramentas necessárias para que sejam bem-sucedidos nas diferentes situações de sua vida cidadã – profissional, cultural e política.
Essa é uma contribuição efetiva que os resultados do Enade podem oferecer às instituições de ensino superior do país.
Um segundo aspecto que ficou evidenciado a partir da análise das respostas foi o pouco conhecimento dos graduandos a respeito de temas exteriores ao âmbito específico de sua profissão, ligados à realidade nacional e global e a outras áreas do conhecimento, objetivo das questões de Formação Geral do Enade.
No exame de 2019, isso ficou evidenciado na diferença de desempenho entre estudantes de áreas de Engenharia e de outras áreas de formação na Questão 1, que tratava de desastres ambientais.
Mesmo sem ter acesso à sua área de formação, ao se analisarem as respostas, evidenciou -se o domínio de um vocabulário especializado por uma parte dos graduandos, enquanto a outra parte limitou -se ao emprego de um vocabulário próprio do senso comum, com palavras genéricas e vagas.
Por exemplo, se o tema tivesse sido “ pandemia ”, os estudantes das áreas da saúde teriam, provavelmente, tido um desempenho vocabular distinto, utilizando termos técnicos desconhecidos pelos estudantes de Engenharia.
Outro ponto a ser observado, já em consequência da análise da Questão 2, é o baixo nível de conhecimento e de reflexão dos estudantes sobre um aspecto tão importante da vida nacional : a relação entre o público e o privado, no que diz respeito ao ensino superior, à pesquisa universitária e seu papel no desenvolvimento brasileiro.
A respeito da adequação do exame aos seus propósitos institucionais, vale observar que a Questão 2 não atendeu às exigências de avaliação, tanto do domínio discursivo dos estudantes, quanto do seu conhecimento sobre temas diversos ligados à realidade nacional e à global, assim como a outras áreas do conhecimento exteriores ao âmbito específico de sua profissão.
Conforme foi detalhadamente explicado no corpo deste relatório, o enunciado da Questão 2 induziu os estudantes a criarem respostas vagas e insatisfatórias, com repertório lexical pouco diversificado, impreciso ou sem desenvolvimento, com tessitura textual comprometida por truncamentos sintático-semânticos, reflexo da fragmentação do pensamento.
Se o objetivo é avaliar a competência discursiva dos formandos brasileiros, a elaboração de um texto dissertativo certamente ofereceria material relevante para essa análise, por incentivar a reflexão e a organização do pensamento em uma sequência textual coerente e complexa. 6.3.2 Componente de Conhecimento Específico Na parte da prova relativa às questões discursivas do Componente de Conhecimento Específico ( Tabela 6.18 ), observa -se que a Média foi mais alta do que para as questões discursivas do Componente de Formação Geral.
Enquanto no Componente de Formação Geral a Média para estudantes de Engenharia Florestal de todo o Brasil foi 34,0, na parte de Conhecimento Específico, a Média foi 36,6.
A maior Média deste componente foi obtida pelos estudantes da região Sudeste ( 41,2 ), e a menor, pelos da região Norte ( 32,0 ).
As demais regiões obtiveram Média igual a 33,4 na região Nordeste, a 39,7, na Sul, e a 36,8 na região Centro-Oeste.
Quanto à variabilidade das notas, o Desvio padrão de todo o Brasil foi 16,9.
O maior Desvio padrão foi encontrado na região Sul ( 17,4 ), e o menor, na região Centro-Oeste ( 15,6 ).
A maior nota Máxima, 85,0, foi obtida na região Sul.
Nas demais regiões, as notas Máximas foram : 81,7, na região Norte, 73,3, na Nordeste, 83,3, na Sudeste e 75,0, na região Centro-Oeste.
A nota Mínima ( 0,0 ) foi obtida por, pelo menos, um estudante em todas as regiões do Brasil.
A Mediana do Brasil foi 36,7, a mesma da região Centro-Oeste.
As demais regiões apresentaram as seguintes Medianas : 31,7, na região Norte, 33,3, na Nordeste, 41,7, na Sudeste, e 39,2, na região Sul.
No Gráfico 6.11, representa -se a distribuição das notas nas questões discursivas no Componente de Conhecimento Específico.
A moda dessa distribuição ocorre no intervalo ( 30 ; 40 ], com 22,9 \% do total de participantes.
A análise de \textbf{cada} uma destas questões será feita a seguir. 6.3.2.1 Análise da Questão Discursiva 3 do Componente de Conhecimento Específico Na questão 3, cujos resultados aferidos se encontram descritos na Tabela 6.19, a Média dos estudantes de todo o Brasil foi 26,5.
O desempenho dos estudantes nessa questão foi inferior ao verificado na questão 4 e superior ao verificado na questão 5.
A menor Média nessa questão foi obtida pelos estudantes da região Nordeste ( 22,6 ), enquanto a maior Média foi obtida na região Norte ( 28,9 ).
Quanto à variabilidade das notas, o Desvio padrão de todo o Brasil foi 23,5.
O maior Desvio padrão foi obtido na região Sudeste ( 24,4 ), enquanto o menor foi obtido na região Centro-Oeste ( 21,5 ).
A nota Máxima, 100,0 pontos, foi alcançada por, pelo menos, um estudante em três regiões, com exceção das regiões Norte e Centro-Oeste cuja nota Máxima foi 90,0, em ambas.
A Mediana do Brasil, como um todo, foi 30,0, a mesma das regiões Norte, Sudeste e Sul.
As regiões Nordeste e Centro-Oeste obtiveram Medianas 15,0 e 20,0, respectivamente.
A nota Mínima ( 0,0 ) foi a mesma em todas as regiões do Brasil.
No Gráfico 6.12, mostra -se a distribuição das notas na questão discursiva 3, do Componente de Conhecimento Específico da área de Engenharia Florestal.
Observa -se que 4,4 \% dos estudantes deixaram a questão em branco e que 26,6 \% receberam nota zero ( classe modal ).
Destacam -se, também, os intervalos ( 0 ; 10 ] e ( 20 ; 30 ] como dois máximos locais, com aproximadamente o mesmo percentual de participantes, 17,4 \% e 17,3 \%, respectivamente. 6.3.2.2 - Comentários sobre as respostas à Questão Discursiva 3 A questão era vinculada ao núcleo de conhecimento estabelecidos pela Diretriz Curricular publicada na Resolução nº 3, de 2 de fevereiro de 2006, para a formação do engenheiro florestal.
Abordava o manejo florestal, conteúdo essencial para a formação profissional.
Na questão exigia -se conhecimento de manejo florestal madeireiro e não madeireiro.
O conhecimento avaliado pela questão pode ser considerado básico e fundamental, sendo o mínimo esperado para um engenheiro florestal.
O padrão de respostas era claro e previa boa variedade de respostas possíveis.
Porém, o item ‘ a ’, que solicitava que o estudante descrevesse duas etapas, gerou muitas possibilidades de respostas, algumas não previstas no padrão estabelecido.
Por exemplo, como uma das opções de resposta para o item ‘ a ’ o padrão de resposta apresentava a seguinte descrição : “ Mapeamento temático : representação das informações geográficas inerentes à área e cruciais para planejamento da logística e definição de cronograma de exploração ”.
No entanto, essa etapa elementar, na prática, envolve a separação de Unidades de Produção Anual ( UPAs ), Unidades de Trabalho ( UTs ), mas não necessariamente a definição de cronograma de exploração, que poderia \textbf{estar} em uma etapa elementar sobre “ Planejamento da exploração florestal ou colheita ”, que não consta no padrão.
Em contrapartida, resposta como “ Levantamento socioeconômico : diagnóstico da área, visando conhecer as comunidades impactadas pelo empreendimento, tanto na área explorada como no seu entorno ”, presente no padrão de resposta, possui explicação que não complementa a citação do tópico, possuindo praticamente a mesma abordagem.
O mesmo pode ser observado para “ Estudo de Impacto Ambiental ”.
Assim, estudantes que optaram por citar essas duas etapas possuíam uma probabilidade muito maior de tirar boas notas, já que a explicação é muito similar à etapa em si.
Já “ Inventário florestal ” e “ Mapeamento temático ” exigiam um nível de \textbf{detalhamento} maior, havendo menor possibilidade de os estudantes detalharem todos eles.
O enunciado da questão \textbf{estava} claro, e ela pode ser considerada de média a fácil resolução em relação às demais questões discursivas de Conhecimentos Específicos, pois deixava livre ao estudante abordar diferentes etapas elementares possíveis do plano de manejo.
Uma vez que descrever uma etapa deixa a possibilidade de resposta ampla, como já comentado, foram observadas algumas respostas não previstas no padrão, conforme mostra -se a seguir.
Respostas corretas, que não constavam no padrão : Item ‘ a ’ “ Planejamento da exploração florestal ” ou “ Planejamento da colheita ”. “ Adequação a legislação pertinente ” ( como : submissão de Plano Operacional Anual – POA, pedido de Autorização de Exploração – AUTEX ) “ Estudo da viabilidade econômica ” ( que não é exigido pelo órgão ambiental competente, mas que, sem lucro, não seria realizado.
Então, ao menos o engenheiro florestal apresentará ao proprietário da área ). “ Mapeamento temático ” ( \textbf{está} no padrão ), mas dificilmente o estudante abordaria a definição de cronograma de exploração, nesta etapa.
Deste modo, a explicação não ocasionava pontuação, visto que ela \textbf{estava} incompleta em relação ao padrão de respostas.
Não significa, no entanto, que o discente tenha respondido de forma errada, mas apenas incompleta quanto ao exigido pelo padrão. “ Inventário Florestal ” ( \textbf{está} no padrão ), mas com relação à explicação poderiam ser considerados “ levantamento florístico e fitossociológico ” ( que se refere à avaliação de estruturas horizontal, vertical, paramétrica e diamétrica, que são formas de quantificação exigidas para a aprovação de planos de manejo ).
Ou ainda, citando levantamento florístico e fitossociológico, mas não utilizando a palavra inventário florestal ( o que realmente é básico ), e que, por igual motivo, não recebia nenhuma pontuação.
Neste sentido, o estudante tem noção do que deve ser feito, porém, não utilizou os termos técnicos adequados.
Item ‘ b ’ “ Dificuldade para determinar o limite de exploração do produto de modo a garantir a sustentabilidade ”, que poderia \textbf{estar} relacionada à “ Dificuldade de quantificar os impactos na floresta ” ( presente no padrão de resposta, mas, da forma como \textbf{está} no padrão, não quer dizer necessariamente a mesma coisa, pois dá uma ideia de posterior à obtenção do produto, e não anterior para definir os limites adequados ).
Respostas incorretas, mas frequentemente encontradas : Item ‘ a ’ Estudo da área como : Análise de solo, condições meteorológicas e climáticas.
Item ‘ b ’ Falta de estudos científicos ; falta de mão de obra especializada ; falta de incentivo do governo ; dificuldade de acesso à área a ser manejada ; alto custo para obtenção do produto e dificuldade de se adequar à legislação ( que na verdade não existe ou não é clara ).
Em relação à linguagem adotada pelos estudantes, a falta de utilização de termos científicos adequados, por vezes, dificultou o entendimento da resposta.
Alguns termos, quando utilizados de forma equivocada, geravam ambiguidade na resposta.
No Item ‘ a ’, a abordagem mais frequentemente encontrada foi a descrição/realização de inventário florestal.
Os principais erros \textbf{estavam} relacionados à descrição das etapas, com muitas informações desnecessárias, não se atendo ao que de fato é o diferencial da referida etapa.
As soluções apresentadas foram basicamente as contidas no padrão, apesar de em muitos casos serem apresentadas de forma incompleta ou equivocada, além de apresentarem como soluções : Planejamento da exploração florestal ou colheita ; Adequação à legislação pertinente e Estudo da viabilidade econômica, sendo estas não encontradas no padrão de respostas.
No item ‘ b ’, as abordagens mais frequentemente encontradas foram a dificuldade de realização de estudos e a falta de mão de obra adequada ( ambas incorretas ) e a dificuldade de se adequarem a questão as questões legais para o manejo de produtos não madeireiros ( todas incorretas ).
Os principais erros \textbf{estavam} relacionados à dificuldade de se adequarem às questões legais para o manejo de produtos não madeireiros.
No geral, o desempenho dos estudantes na Questão 3 pode ser considerado fraco.
De 1.858 participantes, 82 deixaram a resposta a essa questão em branco, outros 27 tiveram suas respostas consideradas como nulas, desconsideradas ou, ainda, eram respostas sobre tema de outra questão da prova ( resposta trocada ).
As demais respostas foram classificadas em três grupos : fracas, medianas e boas, como discriminado a seguir.
As respostas válidas e classificadas como fracas, por receberem notas de zero a 30, constituíram o maior grupo, 1.121 respostas.
Tais notas foram causadas por vários motivos, dentre os quais, respostas curtas, apenas em tópicos, sem tentativa de explicar as etapas no item ‘ a ’.
Respostas com notas de 31 a 70 totalizaram 569.
Já o grupo menos expressivo foi o das 59 respostas que receberam notas maiores do que 70 e foram consideradas boas.
Apenas seis estudantes obtiveram a nota máxima 100.
Adicionalmente, foi possível perceber que uma grande parte dos estudantes não possui, ou possui pouco conhecimento fundamental básico para a realização de um plano de manejo em florestas nativas, tanto madeireiro, quanto não madeireiro.
Em relação aos planos de manejo não madeireiros, a maioria dos estudantes se restringiu a dizer que a maior dificuldade é a falta de estudos, sem especificar o que se pretende com a realização dos estudos.
Os discentes foram muito lacônicos em suas respostas, o que deixava o corretor imaginando o que o discente queria dizer com aquela afirmativa resumida e se a afirmativa \textbf{estava} ligada a uma das alternativas do padrão de resposta.
Além disso, o fato de o tema de manejo não madeireiro não ter metodologia e normas bem consolidadas, pode ter gerado dificuldades e a grande diversidade de respostas encontradas. 6.3.2.3 - Análise da Questão Discursiva 4 do Componente de Conhecimento Específico Na Tabela 6.20, constam as informações relativas à questão 4 do conjunto de questões do Componente de Conhecimento Específico.
O desempenho dos estudantes nessa questão foi superior ao verificado nas questões 3 e 5.
A Média geral do Brasil foi 60,7, sendo a maior Média registrada na região Sudeste ( 69,1 ), e a menor, na região Norte ( 50,0 ).
A nota Máxima ( 100,0 ) foi atingida por, pelo menos, um concluinte de \textbf{cada} região.
A Mediana do Brasil, como um todo, foi 60,0, a mesma para a região Centro-Oeste.
Nas demais regiões, as Medianas foram : 50,0 na região Norte, 55,0, na Nordeste, 75,0, na Sudeste, e 70,0, na região Sul.
A nota Mínima foi zero em todas as regiões.
No Gráfico 6.13, \textbf{está} representada a distribuição de notas da questão discursiva 4, do Componente de Conhecimento Específico.
Observa -se que pouco menos de 5 \% dos participantes deixaram em branco a resposta a essa questão.
O intervalo modal é ( 40 ; 50 ], representado por 18,4 \% das notas, seguido dos intervalos ( 90 ; 100 ], ( 30 ; 40 ] e ( 70 ; 80 ] com cerca de 14 \% de participantes \textbf{cada} intervalo. 6.3.2.4 Comentários sobre as respostas à Questão Discursiva 4 A Questão 4 \textbf{estava} vinculada ao núcleo de conhecimentos estabelecidos pelas Diretrizes Curriculares ( Resolução nº 3, de 2 de fevereiro de 2006 ) de formação do engenheiro florestal.
O tema abordado era Silvicultura, que faz parte do núcleo essencial da área.
O nível de complexidade exigido na Questão 4 pode ser considerado de médio a fácil em relação aos dois itens, ‘ a ’ e ‘ b ’, solicitados.
No item ‘ a ’, solicitou -se a resolução de um cálculo, que teve enunciado claro, com informações suficientes.
Esse cálculo é algo que \textbf{está} contemplado nas competências e habilidades exigidas para um engenheiro florestal.
O padrão de resposta do item apresentava a resolução numa sequência de cálculos, que, de forma geral, são de grau de complexidade fácil a médio.
Apesar de haver espaço suficiente para que os cálculos pudessem ser apresentados, o comando não solicitava ao estudante a apresentação dos cálculos.
Devido à falta de exigência explicita de apresentar os cálculos, houve uma diversidade de situações.
Houve apresentação apenas dos resultados finais dos cálculos, houve os que apresentaram alguns dos resultados intermediários, outros apresentaram todas as etapas No item ‘ a ’, solicitou -se a citação de fatores que possam ser causa de efeito ( morte de mudas ).
O comando também \textbf{estava} claro, com dados suficientes e dentro das competências do engenheiro florestal.
Nesse item não foi solicitado ao estudante justificar/explicar os fatores citados.
Houve boa relação entre a resposta dos estudantes e o padrão de resposta.
Na correção do item ‘ b ’, observou -se que há estudantes que dominam os termos técnicos específicos da área e outros não.
Alguns compreenderam e usaram termos equivalentes ou dissertaram para explicar.
Como no item ‘ b ’ era solicitada apenas a citação, muitas respostas foram sucintas, com uso de termos diferentes do padrão, que ora \textbf{estavam} relacionados a algum item do padrão, ora não foram incorretos ou muito genéricos.
Por exemplo : “ solo ”, embora um solo mal preparado, um solo compactado, um solo com excesso de umidade, etc., possa ser considerado como uma resposta correta para o item ‘ b ’, o estudante ao colocar apenas a palavra “ solo ”, sem dar qualquer complementação, não recebeu pontuação.
Isso se repetiu para vários outros termos.
Também houve respostas dadas por estudantes que apesar de não previstas pelo padrão de respostas, eram corretas e foram consideradas, como por exemplo, “ ocorrência de queimadas ”.
Quanto ao desempenho dos estudantes no item ‘ a ’, houve significativo número de respostas em branco e uma grande quantidade de respostas que apresentavam apenas o resultado final, ora corretos ( com pontuação total ) ora incorretos ( com pontuação zerada ).
Houve também várias respostas com apresentação dos cálculos e dos resultados intermediários, o que permitiu estabelecer pontuação parcial.
Observou -se uma quantidade grande de erros de cálculos ( divisão, multiplicação, adição ).
Também houve, em quantidade pequena, respostas completas, em que os estudantes explicavam detalhadamente a resolução dos cálculos.
No item ‘ b ’, raramente se observaram respostas em branco e com pontuação zerada.
Houve diversidade nas respostas, mas alguns fatores foram frequentemente citados, como por exemplo : “ ataque de pragas ( formigas e cupins ) ”, “ doenças ”, “ plantio mal executado ”, “ mudas não rustificadas ”, dentre outras.
O item ‘ b ’ da questão foi o de melhor desempenho dos estudantes.
Raramente houve uma resposta que não tenha obtido pontuação nesse item, no entanto, houve grande diversidade de respostas com uso de termos às vezes equivalente aos do padrão e ainda respostas com justificativa/explicação.
No geral, o desempenho dos estudantes na Questão 4 pode ser considerado regular.
Do total de 1.858 participantes, 79 deixaram a resposta a essa questão totalmente em branco, outros 14 tiveram suas respostas consideradas como nulas, desconsideradas ou, ainda, eram respostas sobre tema de outra questão da prova ( resposta trocada ).
As demais respostas foram classificadas em três grupos : fracas, medianas e boas, como discriminado a seguir.
As respostas classificadas como fracas, por receberem notas de zero a 30, constituíram o menor grupo, 177 respostas.
Os motivos para essas notas foram variados, mas o que se observou com maior frequência foram respostas em branco para esse item ou a apresentação apenas do valor final errado.
Além disso, no item ‘ ’ b, as respostas eram sucintas, em tópicos, sem justificativa/explicação, e com uso de termos fora do padrão de resposta.
A quantidade de respostas com notas de 31 a 70 foi 872, aproximadamente a metade das respostas válidas corrigidas.
Essas respostas normalmente apresentavam os cálculos feitos pelo estudante e, dessa forma, foi possível pontuar parcial ou integralmente o item ‘ a ’.
Nessa faixa de notas, também foi frequente o item ‘ a ’ receber nota zero por \textbf{estar} errado, por erros dos cálculos e do valor final ou por não apresentar resposta ( em branco ).
Já no item ‘ b ’, as respostas, apesar de sucintas e em tópicos, foram feitas de forma correta ou parcialmente correta dentro do padrão.
Houve com frequência justificativa/explicação dos termos que \textbf{estavam} fora do padrão de resposta, mas que equivaliam à algum termo do padrão.
Já as 716 respostas que receberam notas maiores do que 70, consideradas boas, apresentavam os cálculos feitos pelo estudante e, dessa forma, foi possível pontuar parcial ou integralmente o item ‘ a ’, ou então foi apresentado apenas o valor final correto.
Já no item ‘ b ’, apesar de respostas sucintas e em tópicos, eram corretas ou parcialmente corretas.
Com frequência, nas respostas boas, havia justificativa/explicação para os fatores citados.
Nesse grupo, 258 estudantes obtiveram a nota máxima 100.
A questão, de forma geral, era básica para um engenheiro florestal, exigindo -se uma capacidade de resolução de um cálculo ( não complexo ) e a análise de fatores causadores de morte de mudas.
Pareceu haver uma dificuldade na execução dos cálculos, contas relativamente simples, mas que, sem calculadora, uma grande quantidade de estudantes não conseguiu executar de forma correta, o que é preocupante para uma formação em engenharia.
Por não terem sido exigidos os cálculos intermediários, não foi possível ter uma melhor compreensão da dificuldade do estudante na estruturação da resolução do solicitado.
Se ele foi capaz de entender o que foi pedido e errou apenas na execução do cálculo, ou se não conseguiu compreender o que foi pedido ou se usou formas erradas de encontrar o resultado.
Observou -se grande quantidade de respostas muito sucintas, o que pode ter sido Na Tabela 6.21, constam as informações relativas à questão 5 do conjunto do Componente de Conhecimento Específico.
O desempenho dos estudantes de todo o Brasil nesta questão foi o mais baixo entre as três questões discursivas desse componente.
A nota Média dos estudantes de todo o Brasil foi 22,4.
A maior Média foi registrada nas regiões Sudeste e Sul, ambas, com 26,3, enquanto a menor Média foi registrada na região Norte ( 17,1 ).
Quanto à variabilidade das notas, o Desvio padrão dos estudantes do Brasil, como um todo, foi 21,1.
Enquanto o maior desvio foi encontrado na região Sudeste ( 22,9 ), o menor foi encontrado na região Norte ( 17,6 ).
A Mediana do Brasil, como um todo, foi 20,0, a mesma para as regiões Sul e Centro-Oeste.
Nas demais regiões as Medianas foram 15,0, nas regiões Norte e Nordeste, e 25,0, na região Sudeste.
Para o conjunto de estudantes de Engenharia Florestal do Brasil, a nota Máxima foi 95,0, obtida por, pelo menos, um estudante nas regiões Sudeste e Sul.
As outras três regiões obtiveram nota Máxima 85,0.
A nota Mínima do Brasil foi zero, o mesmo valor encontrado em todas as regiões.
No Gráfico 6.14, é representada a distribuição das notas da questão discursiva 5 do Componente de Conhecimento Específico.
Essa distribuição tem o intervalo moda no intervalo ( 10 ; 20 ], correspondendo a 19,9 \% dos respondentes.
Os participantes que tiraram zero, tendo respondido à questão representam 16,7 \% do total.
Somando -se a frequência desse intervalo com a daqueles que deixaram a questão em branco, chega a 26,8 \% os participantes que zeraram a questão. devido ao que à forma como os comandos da questão estabeleceram a tarefa.
Pode também ser fruto do pelo perfil dos estudantes. 6.3.2.5 Análise da Questão Discursiva 5 do Componente de Conhecimento Específico Na Tabela 6.21, constam as informações relativas à questão 5 do conjunto do Componente de Conhecimento Específico.
O desempenho dos estudantes de todo o Brasil nesta questão foi o mais baixo entre as três questões discursivas desse componente.
A nota Média dos estudantes de todo o Brasil foi 22,4.
A maior Média foi registrada nas regiões Sudeste e Sul, ambas, com 26,3, enquanto a menor Média foi registrada na região Norte ( 17,1 ).
Quanto à variabilidade das notas, o Desvio padrão dos estudantes do Brasil, como um todo, foi 21,1.
Enquanto o maior desvio foi encontrado na região Sudeste ( 22,9 ), o menor foi encontrado na região Norte ( 17,6 ).
A Mediana do Brasil, como um todo, foi 20,0, a mesma para as regiões Sul e Centro-Oeste.
Nas demais regiões as Medianas foram 15,0, nas regiões Norte e Nordeste, e 25,0, na região Sudeste.
Para o conjunto de estudantes de Engenharia Florestal do Brasil, a nota Máxima foi 95,0, obtida por, pelo menos, um estudante nas regiões Sudeste e Sul.
As outras três regiões obtiveram nota Máxima 85,0.
A nota Mínima do Brasil foi zero, o mesmo valor encontrado em todas as regiões.
No Gráfico 6.14, é representada a distribuição das notas da questão discursiva 5 do Componente de Conhecimento Específico.
Essa distribuição tem o intervalo moda no intervalo ( 10 ; 20 ], correspondendo a 19,9 \% dos respondentes.
Os participantes que tiraram zero, tendo respondido à questão representam 16,7 \% do total.
Somando -se a frequência desse intervalo com a daqueles que deixaram a questão em branco, chega a 26,8 \% os participantes que zeraram a questão. 6.3.2.6 Comentários sobre as respostas à Questão Discursiva 5 A Questão 5 \textbf{estava} vinculada ao núcleo de conhecimento estabelecidos pelas Diretrizes Curriculares ( Resolução nº 3, de 2 de fevereiro de 2006 ).
O conteúdo avaliado foi Tecnologia e Utilização dos Produtos Florestais.
A questão exigia conhecimento de tecnologia da madeira, relação água madeira, noções de anatomia da madeira e como a densidade da madeira se vincula à secagem, processo em que a madeira perde sua água natural.
No item ‘ a ’, foi pedida a explicação da relação entre a densidade básica da madeira e sua secagem.
Além de ser uma questão aberta, o comando não deixava claro que tipo de abordagem se esperava do discente.
A resposta correta mais frequente foi que a densidade básica influencia no tempo de secagem, seja natural ou em estufa.
Poucas foram as respostas em que os estudantes conectaram a densidade básica à quantidade de fração parede ou porosidade, ou mesmo quanto à capacidade de retenção de água, seja na parede celular ou nos vazios celulares.
O comando poderia ter deixado claro que se esperava a relação entre densidade básica e a velocidade de secagem, o teor de umidade da madeira, sua porosidade e defeitos oriundos da secagem.
Assim, evitar -se-ia a grande quantidade de abordagens distintas encontrada na correção.
No item ‘ b ’ da questão, foi solicitado que o discente citasse três fatores inerentes ao processo de secagem e como eles influenciam na perda de água da madeira.
A indagação \textbf{estava} vinculada ao conhecimento sobre secagem da madeira, que faz parte do núcleo de conhecimentos essenciais da formação de engenheiro florestal.
Os fatores inerentes, seja na secagem ao ar ou em estufa, são a temperatura e a umidade relativa do ambiente e também depende do fluxo ou movimentação do ar na pilha de secagem.
A correção do item ‘ b ’ da questão, deixou claro que a não compreensão da palavra inerente pode ter levado a respostas diversas sobre os fatores que influenciam a secagem, devido a interpretações diferentes da esperada pelo padrão de resposta.
O nível de exigência de conhecimentos sobre secagem de madeira da questão pode ser considerado básico.
No entanto, percebeu -se dificuldade de desenvolvimento das respostas.
No item ‘ a ’ da questão, ficou evidente que o conceito de densidade básica, conceito elementar em tecnologia de produtos florestais, foi usado de forma equivocada por muitos daqueles que fizeram a prova.
Apesar do enunciado não pedir definição de densidade básica, foi muito comum encontrar respostas com definição do que era densidade básica/densidade aparente e a definição de secagem da madeira.
Muitos não conseguiram associar, ou relacionar, a secagem da madeira com a densidade básica.
Alguns outros confundiram densidade básica com densidade aparente, afirmando que a densidade básica se altera à medida que o teor de umidade da madeira foi alterado no processo de secagem.
Poucos discentes conseguiram escrever sobre defeitos de secagem associados à densidade básica da madeira, apesar de o enunciado da questão fazer menção aos defeitos oriundos no processo de secagem.
Alguns poucos discentes chegaram perto da totalidade de pontos do item ‘ a ’ ( 40 pontos ), conseguindo associar a maior densidade básica ao maior tempo de secagem, quando comparado a madeiras de menor densidade, assim como associar a densidade à porosidade da madeira e dissertando como a densidade básica se relaciona com o teor de umidade da madeira.
Alguns, conseguiram distinguir os tipos de águas existentes na madeira e as distintas dificuldades em retirá -las.
No item ‘ a ’, a resposta correta de maior ocorrência foi a associação entre densidade básica e velocidade de secagem.
Poucas vezes apareceu a associação entre densidade e defeitos de secagem.
A abordagem equivocada mais comum foi a definição de densidade básica ou densidade aparente, e posteriormente a tentativa de associação com a secagem e a umidade da madeira.
Responder que o processo de secagem afetava a densidade básica da madeira, interferindo em sua massa ou/e seu volume ocorreu várias vezes.
Alguns meramente definiram o que é densidade básica, muitas vezes de forma equivocada, e buscando associar a densidade à secagem, mas sem sucesso.
A abordagem correta que mais apareceu como resposta foi a relação entre a densidade básica e a velocidade de secagem ( associada ao item ‘ b ’ ).
No item ‘ b ’ da questão, os fatores inerentes ao processo de secagem são ambientais ( natural ou artificial ), envolvendo temperatura, umidade relativa do ar e circulação ou movimentação do ar.
Como já foi citado anteriormente, a possibilidade da não compreensão da palavra inerente pode ter gerado a multiplicidade de respostas encontrada na correção.
Os fatores inerentes foram confundidos com técnicas usadas no processo de secagem para melhorá -la ou mesmo para se evitarem maiores riscos de defeitos.
Assim, foram citados : cuidado na formação das pilhas de secagem com adequado entabicamento ; secagem de madeira da mesma espécie e mesma espessura ; posicionamento da pilha ( expressão muito vaga ) ; como montar ou construir uma pilha de secagem ; não colocar a madeira em temperatura muito alta no início do processo de secagem, além da descrição dos cuidados na organização e montagem de um pátio de secagem.
Também apareceram respostas envolvendo questões anatômicas, os aspectos da constituição química da madeira ( teor de lignina, presença de extrativos, resinas etc. ), com tentativas de explicar como esses fatores, anatômicos e químicos, afetavam o processo de secagem.
Foi citado muitas vezes que a densidade básica da madeira e sua porosidade eram fatores intrínsecos.
Por vezes, quando o discente citava apenas a palavra umidade não ficava claro a qual umidade ele se referia, se era à umidade da madeira ou à umidade relativa do ar.
Nesse caso, a resposta foi considerada \textbf{certa} ou errada, após a leitura da explicação.
No item ‘ b ’, houve uma grande diversidade de respostas.
A abordagem equivocada mais presente foi a de citar as técnicas de secagem que procuram garantir uma secagem mais eficiente, técnicas tais como as de : adequado empilhamento da madeira ; construção da pilha, posicionamento da pilha num pátio de secagem em ambiente natural ou mesmo em estufa ; padronização das dimensões da madeira a ser seca ; separação por espécie.
Muitos discentes citaram a densidade como fator usado na tomada de decisão de como a madeira deveria ser seca, se em estufa ou no ambiente natural e, quando secas em estufa, como deveria ser a escolha do programa de secagem, procurando assim reduzir o tempo de secagem ou mesmo os defeitos resultantes do processo.
Muito do que foi escrito era verdadeiro quanto ao processo adequado de secagem, mas fugia aos fatores inerentes estabelecidos no padrão de resposta.
Tempo de secagem apareceu com fator inerente em muitas das respostas, de forma equivocada.
O tempo de secagem é decorrente, entre outros fatores, dos fatores inerentes apresentados no padrão de resposta.
A abordagem correta mais frequentemente utilizada foi o fator temperatura do ambiente e a umidade relativa do ar.
Quanto à explicação dos fatores, nem sempre a explicação foi satisfatória para que o discente recebesse pontuação.
No geral, o desempenho dos estudantes na Questão 5 pode ser considerado fraco, com média 22,4.
Do total de 1.858 participantes, 188 deixaram a resposta a essa questão em branco, outros seis tiveram suas respostas consideradas como nulas, desconsideradas ou, ainda, eram respostas sobre tema de outra questão da prova ( resposta trocada ).
As demais respostas foram classificadas em três grupos : fracas, medianas e boas, como discriminado a seguir.
As respostas classificadas como fracas, por receberem notas de zero a 30, constituíram o maior grupo, 1.139 respostas.
Nesse grupo predominaram respostas curtas, apenas em tópicos, sem tentativa de explicar ou discorrer sobre o tema relação densidade básica e secagem, nem sobre os fatores inerentes ao processo de secagem.
Quando houve tentativa de relacionar, muitas vezes as explicações eram superficiais, não correspondendo ao esperado pelo padrão, ou então, eram mesmo equivocadas.
A quantidade de respostas com notas de 31 a 70, consideradas medianas, foi 460.
As respostas medianas foram aquelas em que, ao menos, os discentes tentaram discorrer sobre a relação densidade e secagem e conseguiram identificar ao menos um dos fatores inerentes ao processo de secagem ou ainda explicaram como esse fator afeta a secagem.
A maioria das notas medianas acima de 50 pontos resultaram de acerto no item ‘ b ’.
Nas 65 respostas boas, com notas acima de 70, os discentes tiveram melhor desempenho no item ‘ b ’ da questão.
No item ‘ a ’, a explicação sobre densidade e velocidade de secagem apareceu predominantemente em quase todas as respostas desse grupo.
Nas respostas boas, os termos técnicos foram usados com mais precisão e frequência, e a escrita era mais clara e objetiva.
Nenhum estudante obteve a nota máxima 100.
Foi possível identificar deficiência de expressão dos conhecimentos e dificuldade com o uso de termos técnicos, visto que, em muitos casos, surgiam dúvidas sobre o que realmente o discente pretendia responder, dificultando a correção.
Ficou evidente que no item ‘ a ’ da questão o desempenho dos discentes foi inferior ao verificado no item ‘ b ’.
Os conhecimentos sobre os fundamentos de tecnologia da madeira ( densidade básica, tipos de águas existentes, parede celular e vazios celulares, capacidade de retenção de água da parede celular ou nos vazios ) mostraram -se confusos.
Foi demonstrado um melhor conhecimento sobre o processo de secagem da madeira ( item ‘ b ’ ), entre aqueles que fizeram a prova.
Houve textos longos com notas baixas podendo -se aproveitar muito pouco daquilo que foi escrito.
Muitas vezes, as respostas mais objetivas obtiveram as melhoras notas, mostrando que aqueles que tinham conhecimento sobre o tema também tinham capacidade de síntese.nnA Questão 5 foi aquela que teve menor desempenho, quando comparada às outras duas questões discursivas de Conhecimentos Específicos.
A quantidade de respostas em branco foi muito superior à das demais questões discursivas.nn6.3.2.7 Considerações Finais sobre as questões discursivas do Componente de Conhecimentos EspecíficosnAs questões discursivas da prova de Engenharia Florestal \textbf{estavam} relacionadas ao núcleo de formação essencial da formação do engenheiro florestal, estabelecido pelas Diretrizes Curriculares.
Os conteúdos avaliados foram Manejo Florestal, Silvicultura e Tecnologia e Utilização dos Produtos Florestais.
As questões não apresentaram alto nível de dificuldade, mas as médias das notas de duas delas foram baixas.
A Questão 4, sobre Silvicultura, que envolvia os cálculos de produção de mudas para implantar uma determinada área e que continha também os motivos que poderiam levar à necessidade de replantio de mudas alcançou a maior média.
No entanto, observou -se que os discentes tiveram dificuldades com cálculos relativamente simples ( cálculo de área, multiplicação, divisão e cálculo de percentagem ), erros que fizeram com que as notas do item ‘ a ’ fossem menores que as do item ‘ b’.nnA questão de manejo florestal madeireiro e não madeireiro obteve média final baixa, revelando fraco desempenho por parte dos estudantes.
Ficou claro que uma parte dos discentes não possuía ou possuía pouco conhecimento básico para a realização de um plano de manejo madeireiro e não madeireiro.
Tal resultado pode \textbf{estar} relacionado a fatores, tais como : as dimensões do Brasil, com vários biomas diferentes e distintas realidades florestais, e a \textbf{localização} da Universidade formadora, que influiria na apresentação de conteúdo de manejo florestal de nativas para fins madeireiro e não madeireiro.nnQuanto à questão de tecnologia e utilização de produtos florestais, observou -se que o desempenho discente quanto ao conhecimento de conteúdo e dos conceitos básicos foram fracos.
Além disso, foi a questão que apresentou o maior número de respostas deixadas em branco.

% matched lemmas: cada, certo, detalhamento, estar, localização, torno
\end{textsample}
